\pdfoutput=1 % ensure arXiv direct PDF output

\documentclass[11pt]{amsart}

\usepackage{amssymb,amsthm,enumitem,colonequals,tikz-cd,microtype}
\usepackage[normalem]{ulem} 

%% Font style: old versions mlmodern, \usepackage[cal=euler,bfcal,bb=px,bfbb]{mathalpha}
\usepackage[osf]{Baskervaldx}
\usepackage[baskervaldx]{newtxmath}
\usepackage[cal=boondoxo]{mathalfa} % mathcal from STIX, unslanted a bit

%% Geometry with custom margins
\usepackage[top=3.75cm, bottom=3cm, left=3.5cm, right=3.5cm]{geometry}

%% High penalty to avoid footnote breaks
\interfootnotelinepenalty=10000 

%% Color definitions for links and references
\usepackage{xcolor}
\colorlet{darkblue}{blue!55!black}
\colorlet{darkcyan}{cyan!50!black}
\colorlet{darkgreen}{green!60!black}

%% URL hyphenation option must come before hyperref
\PassOptionsToPackage{hyphens}{url}
\usepackage{hyperref}
\hypersetup{
    colorlinks=true,
    linkcolor=darkblue,
    urlcolor=darkcyan,
    citecolor=darkgreen,
}

%% Colored eqref for emphasis (optional)
\def\eqref#1{\textcolor{darkblue}{(\ref{#1})}}

%% Cleveref utility
\usepackage[nameinlink]{cleveref} 
\Crefformat{section}{#2\S#1#3}
\Crefmultiformat{section}{#2\S\S#1#3}{ and~#2#1#3}{, #2#1#3}{, and~#2#1#3}

%% Overfill notice
\usepackage[pagewise]{lineno}
\overfullrule=100pt % for debugging, comment or set to 0 for final

\let\oldequation\equation
\let\oldendequation\endequation
\renewenvironment{equation}{\linenomathNonumbers\oldequation}{\oldendequation\endlinenomath}
\expandafter\let\expandafter\oldequationstar\csname equation*\endcsname
\expandafter\let\expandafter\oldendequationstar\csname endequation*\endcsname
\renewenvironment{equation*}{\linenomathNonumbers\oldequationstar}{\oldendequationstar\endlinenomath}
\let\oldalign\align
\let\oldendalign\endalign
\renewenvironment{align}{\linenomathNonumbers\oldalign}{\oldendalign\endlinenomath}
\expandafter\let\expandafter\oldalignstar\csname align*\endcsname
\expandafter\let\expandafter\oldendalignstar\csname endalign*\endcsname
\renewenvironment{align*}{\linenomathNonumbers\oldalignstar}{\oldendalignstar\endlinenomath}

%% Theorem and counter setup
\theoremstyle{plain}
\newtheorem{theorem}{Theorem}[section]
\newtheorem{lemma}[theorem]{Lemma}
\newtheorem{corollary}[theorem]{Corollary}
\newtheorem{proposition}[theorem]{Proposition}

\theoremstyle{definition}
\newtheorem{definition}[theorem]{Definition}
\newtheorem{example}[theorem]{Example}
\newtheorem{remark}[theorem]{Remark}
\newtheorem{setup}[theorem]{Setup}
\newtheorem{notation}[theorem]{Notation}

\newtheorem*{disclaimer}{Disclaimer}
\newtheorem*{ack}{Acknowledgments}

% declare alias for cleveref, otherwise it doesn't do what it should; this probably isn't best fix but found it on the web
% \AddToHook{env/introthm/begin}{\crefalias{intro}{theorem}}
\AddToHook{env/notation/begin}{\crefalias{theorem}{notation}}
\AddToHook{env/lemma/begin}{\crefalias{theorem}{lemma}}
\AddToHook{env/corollary/begin}{\crefalias{theorem}{corollary}}
\AddToHook{env/proposition/begin}{\crefalias{theorem}{proposition}}
\AddToHook{env/definition/begin}{\crefalias{theorem}{definition}}
\AddToHook{env/remark/begin}{\crefalias{theorem}{remark}}
\AddToHook{env/setup/begin}{\crefalias{theorem}{setup}}
\AddToHook{env/example/begin}{\crefalias{theorem}{example}}

%% toc and numbering depths
\setcounter{tocdepth}{1}
\setcounter{secnumdepth}{3}
\numberwithin{equation}{section}
\numberwithin{theorem}{section}

%% Comment out collaborator notes before final version
\newcommand{\Pat}[1]{{\color{blue} $\clubsuit$ Pat: [#1]}}
\newcommand{\Bregje}[1]{{\color{red} $\heartsuit$ Bregje: [#1]}}

\title[Gluing weak approximability along recollements]{Gluing weak approximability along recollements}

\author[P.~Lank]{Pat Lank}
\address{P.~Lank,
Dipartimento di Matematica “F. Enriques”, Universit\`{a} degli Studi di Milano, Via Cesare
Saldini 50, 20133 Milano, Italy}
\email{plankmathematics@gmail.com}

\author[B. ~Pauwels]{Bregje Pauwels}
\address{B. ~Pauwels,
School of Mathematical \& Physical Sciences\\
Macquarie University\\
Sydney, Australia}
\email{bregje.pauwels@mq.edu.au}

\date{\today}

\keywords{Weakly approximable triangulated categories, recollements, $t$-structures, algebraic spaces}

\subjclass[2020]{18G80 (primary), 14F08, 14A20} 

%% 18G80 Derived categories, triangulated categories

%% 14F08 Derived categories of sheaves, dg categories, and related constructions in algebraic geometry 

%% 14A20 Generalizations (algebraic spaces, stacks)

\begin{document}

\begin{abstract}
    We study whether weak approximability glues along recollements of triangulated categories that are compactly generated by a single object.
    This is applied to derived categories of complexes with quasi-coherent cohomology and prescribed support for quasi-compact quasi-separated algebraic spaces.
\end{abstract}

\maketitle

\tableofcontents

%%%%%%%%%%%%%%%%%%%%%%%%%%%%%%%%%%%%%%%%%%%
\section{Introduction}
\label{sec:intro}
%%%%%%%%%%%%%%%%%%%%%%%%%%%%%%%%%%%%%%%%%%%

Approximable triangulated categories, and their weak variant, were introduced by Neeman \cite{Neeman:2021c,Neeman:2021,Neeman:2026}.
The theory has been used to settle conjectures of Bondal--Van den Bergh \cite{Neeman:2021} and of Antieau--Gepner--Heller \cite{Neeman:2024}.
It has also found applications to Fourier--Mukai theory for singular schemes  \cite{Dutta/Lank/ManaliRahul:2025,GuisadoVillaalgordo/Lank/ManaliRahul/Pavic:2025,GuisadoVillaalgordo/Lank:2026}.
See \cite{Canonaco/Neeman/Stellari:2024} for a survey of recent developments.

Roughly speaking, a triangulated category with small coproducts and a compact generator $G$ is weakly approximable if every object in the aisle of a suitable $t$-structure can be suitably approximated from small coproducts of a uniformly bounded range of shifts of $G$.
For example, if $X$ is a quasi-compact quasi-separated scheme and $Z\subseteq X$ is a closed subset with quasi-compact complement, $D_{\operatorname{qc},Z}(X)$ is weakly approximable \cite[Theorem 3.2(iv)]{Neeman:2024}.

The main result is the following. 

\begin{theorem}
    \label{thm:glueing_weak_approx}
    Consider a recollement of singly compactly generated triangulated categories:
    \begin{displaymath}
        \begin{tikzcd}[column sep=6em]
            {\mathcal{D}_Z} & {\mathcal{D}} & {\mathcal{D}_U.}
            \arrow["{i_\ast}" description, from=1-1, to=1-2]
            \arrow["{i^\ast}"', bend right=35, from=1-2, to=1-1]
            \arrow["{i^!}", bend left=35, from=1-2, to=1-1]
            \arrow["{j^\ast}" description, from=1-2, to=1-3]
            \arrow["{j_!}"', bend right=35, from=1-3, to=1-2]
            \arrow["{j_\ast}", bend left=35, from=1-3, to=1-2]
        \end{tikzcd}
    \end{displaymath}
    Assume that $\mathcal{D}_Z$ and $\mathcal{D}_U$ are weakly approximable. Suppose $G^\prime$ is a compact generator of $\mathcal{D}$.
    Then $\mathcal{D}$ is weakly approximable if, and only if, $\operatorname{Hom}(G^\prime,G^\prime[n])=0$ for $n\gg 0$.
\end{theorem}

As an application of \Cref{thm:glueing_weak_approx}, we establish weak approximability for quasi-compact quasi-separated algebraic spaces.

\begin{corollary}
    \label{cor:weak_approx_spaces}
    Let $X$ be a quasi-compact quasi-separated algebraic space with quasi-finite separated diagonal.
    Suppose $S\subseteq |X|$ is a closed subset with quasi-compact complement.
    Then $D_{\operatorname{qc},S}(X)$ is weakly approximable. 
    In particular, $D_{\operatorname{qc}}(X)$ is weakly approximable.
    If $X$ is Noetherian, then $D_{\operatorname{qc},S}(X)$ is coherent.
\end{corollary}

\Cref{thm:glueing_weak_approx} is the analog for weak approximability of \cite[Theorem 5.1]{Burke/Neeman/Pauwels:2023} by Burke--Neeman--Pauwels.
Loc.\ cit.\ shows that $\mathcal{D}$ is approximable if $\mathcal{D}_Z$ and $\mathcal{D}_U$ are approximable and $\operatorname{Hom}(G^\prime,G^\prime[n])=0$ for $n\gg 0$.
The `only if' direction of \Cref{thm:glueing_weak_approx} is formal (cf.\ \cite[Remark 3]{Burke/Neeman/Pauwels:2023}).
Thus, the work lies in the `if' direction.

The proof in \cite{Burke/Neeman/Pauwels:2023} lifts approximations from $\mathcal{D}_Z$ to $\mathcal{D}$.
This is done by induction on the number of cones needed to build them \cite[Lemmas 5.2 to 5.5]{Burke/Neeman/Pauwels:2023}.
At each step, the bound on the lift is worse than on its target.
Hence, the constants grow with the number of cones.
Approximability of $\mathcal{D}_Z$ bounds this number.

No such bound need exist if $\mathcal{D}_Z$ is only weakly approximable.
Instead, we lift a single shift of a compact generator so that it lies in the same shifted aisle as its target (e.g.\ $\mathcal{D}^{\leq K}$).
See \Cref{lem:onecell}.
Consequently, the cone of a lift satisfies the same bound, which gives control on the desired objects.
See \Cref{prop:closure}.

We discuss \Cref{cor:weak_approx_spaces}.
In general for schemes, $D_{\operatorname{qc},Z}$ need not be approximable \cite[Remark 8.1]{Neeman:2024}.
Hence, \cite[Theorem 5.1]{Burke/Neeman/Pauwels:2023} does not apply, and the weak analog is necessary.
The proof of \Cref{cor:weak_approx_spaces} combines Neeman's result for schemes \cite[Theorem 3.2(iv)]{Neeman:2024} with the induction principle for algebraic spaces.

\begin{ack}
    The authors thank Jack Hall and Fei Peng for discussions. 
    Lank was supported by the ERC Advanced Grant 101095900-TriCatApp.
\end{ack}

%%%%%%%%%%%%%%%%%%%%%%%%%%%%%%%%%%%%%
\paragraph*{\textbf{AI disclosure}}
%%%%%%%%%%%%%%%%%%%%%%%%%%%%%%%%%%%%%

Large language models (LLMs) were used to assist with proofreading, improving the exposition, checking arguments, and revising some proofs.

%%%%%%%%%%%%%%%%%%%%%%%%%%%%%%%%%%%%%%%%%%%
\section{Preliminaries}
\label{sec:prelim}
%%%%%%%%%%%%%%%%%%%%%%%%%%%%%%%%%%%%%%%%%%%

%%%%%%%%%%%%%%%%%%%%%%%%%%%%%%%%%%%%%%%%%%%
\subsection{Triangulated categories}
\label{sec:prelim_triangulated}
%%%%%%%%%%%%%%%%%%%%%%%%%%%%%%%%%%%%%%%%%%%

%%%%%%%%%%%%%%%%%%%%%%%%%%%%%%%%%%%%%%%%%%%
\subsubsection{Terminology}
\label{sec:prelim_terminology}
%%%%%%%%%%%%%%%%%%%%%%%%%%%%%%%%%%%%%%%%%%%

We fix some terminology. 
See \cite[\S 1]{Neeman:2001}.
Let $\mathcal{T}$ be a triangulated category.
Note that $\operatorname{Hom}(X,Y)$ is an abelian group.
Given a morphism $u$, write $-u$ for its additive inverse.
Consider a distinguished triangle
\begin{equation}
    \label{eq:triangle}
    X\xrightarrow{u} Y\xrightarrow{v} Z\xrightarrow{w} X[1].
\end{equation}
Then TR2 asserts that
\begin{displaymath}
    \begin{aligned}
        Y
        & \xrightarrow{v} Z\xrightarrow{w} X[1]\xrightarrow{-u[1]} Y[1],
        \\& Z[-1]\xrightarrow{-w[-1]} X\xrightarrow{u} Y\xrightarrow{v} Z
    \end{aligned}
\end{displaymath}
are also distinguished.
Passing from \eqref{eq:triangle} to the first (resp.\ second) of these is called \textbf{rotating} (resp.\ \textbf{rotating backwards}).
Observe that negating two of the three morphisms of a distinguished triangle gives a distinguished triangle. 
Indeed, $(1_X,-1_Y,1_Z)$ is an isomorphism of triangles from $(u,v,w)$ to $(-u,-v,w)$.
% (Conventions in the literature for TR2 differ by such sign changes.)
For every $n\in \mathbb{Z}$, there exists a distinguished triangle
\begin{displaymath}
    X[n]\xrightarrow{u[n]} Y[n]\xrightarrow{v[n]} Z[n]\xrightarrow{(-1)^n w[n]} X[n+1].
\end{displaymath}
We call this \textbf{shifting} the given distinguished triangle by $[n]$.
By TR1, every $f\colon X\to Y$ can be completed into a distinguished triangle 
\begin{displaymath}
    X\xrightarrow{f} Y\to C_f\to X[1].
\end{displaymath}
We call $C_f$ a \textbf{cone} of $f$.
Rotating backwards gives
\begin{displaymath}
    F_f\to X\xrightarrow{f} Y\to F_f[1].
\end{displaymath}
We call $F_f := C_f[-1]$ a \textbf{cocone} of $f$.
Both are unique up to (non-unique) isomorphism.
By combining TR1 and TR2, it follows that every morphism can be realized as the first, second, or third of some distinguished triangle.
Moreover, given a morphism of distinguished triangles, if two of the three vertical morphisms are isomorphisms, then so is the third.
Indeed, this can be shown by rotating and applying \cite[Proposition 1.1.20]{Neeman:2001}.

%%%%%%%%%%%%%%%%%%%%%%%%%%%%%%%%%%%%%%%%%%%
\subsubsection{Generation}
\label{sec:prelim_generation}
%%%%%%%%%%%%%%%%%%%%%%%%%%%%%%%%%%%%%%%%%%%

Let $\mathcal{T}$ be a triangulated category admitting small coproducts and $\mathcal{S},\mathcal{S}^\prime$ full subcategories.
Write $\mathcal{S} \star \mathcal{S}^\prime$ for the full subcategory of objects $X$ admitting a distinguished triangle $S\to X\to S^\prime \to S[1]$ with $S\in \mathcal{S}$ and $S^\prime \in \mathcal{S}^\prime$.
Write $\operatorname{Coprod}(\mathcal{S})$ for the smallest strictly full subcategory containing $\mathcal{S}$ that is closed under small coproducts and extensions, and $\operatorname{smd}(\mathcal{S})$ for the full subcategory of direct summands of objects of $\mathcal{S}$.
For $H\in \mathcal{T}$ and integers $a\leq b$ put
\begin{displaymath}
    H[a,b] := \{ H[-k] : a\leq k\leq b \},
    \qquad 
    \overline{\langle H\rangle}^{[a,b]} := \operatorname{smd}(\operatorname{Coprod}(H[a,b])).
\end{displaymath}

The subcategory of compact objects is denoted by $\mathcal{T}^c$, consisting of those $c \in \mathcal{T}$ for which $\operatorname{Hom}(c,-)$ preserves coproducts.
These form a triangulated subcategory of $\mathcal{T}$. 
The \textbf{localizing subcategory} generated by a collection $\mathcal{S}\subseteq \mathcal{T}$, denoted by $\overline{\langle \mathcal{S} \rangle}^{(-\infty,\infty)}$, is the smallest triangulated subcategory containing $\mathcal{S}$ and closed under small coproducts.
If $\mathcal{T}^c$ is essentially small, then we say that $\mathcal{T}$ is \textbf{compactly generated} when $\mathcal{T} = \overline{\langle \mathcal{T}^c \rangle}^{(-\infty,\infty)}$.
Equivalently, $\mathcal{T}$ is compactly generated if, for any $E \in \mathcal{T}$ satisfying $\operatorname{Hom}(P, E) = 0$ for all $P \in \mathcal{T}^c$, one has $E \cong 0$.
An object $G\in \mathcal{T}^c$ is a \textbf{compact generator} of $\mathcal{T}$ if every $E\in \mathcal{T}$ with $\operatorname{Hom}(G[n],E)=0$ for all $n\in \mathbb{Z}$ satisfies $E\cong 0$.
An object of $\mathcal{T}^c$ is a \textbf{classical generator} of $\mathcal{T}^c$ if the smallest thick subcategory of $\mathcal{T}^c$ containing it is $\mathcal{T}^c$.
If $\mathcal{T}$ is compactly generated, then classical generators for $\mathcal{T}^c$ coincide with compact generators for $\mathcal{T}$ \cite[\href{https://stacks.math.columbia.edu/tag/09SR}{Tag 09SR}]{StacksProject}.
We say that $\mathcal{T}$ is \textbf{singly compactly generated} if it admits a compact generator. 

%%%%%%%%%%%%%%%%%%%%%%%%%%%%%%%%%%%%%%%%%%%%
\subsubsection{\texorpdfstring{$t$}{t}-structures}
\label{sec:prelim_t-structures_t_structures}
%%%%%%%%%%%%%%%%%%%%%%%%%%%%%%%%%%%%%%%%%%%%  

Let $\mathcal{T}$ be a triangulated category with shift functor $[1]$. 
A pair of strictly full subcategories $(\mathcal{T}^{\leq 0}, \mathcal{T}^{\geq 0})$ of $\mathcal{T}$ is a \textbf{$t$-structure} if:
\begin{itemize}
    \item $\operatorname{Hom}(A,B) = 0$ for all $A \in \mathcal{T}^{\leq 0}$ and $B \in \mathcal{T}^{\geq 0}[-1]$,
    \item $\mathcal{T}^{\leq 0}[1] \subseteq \mathcal{T}^{\leq 0}$ and $\mathcal{T}^{\geq 0}[-1] \subseteq \mathcal{T}^{\geq 0}$,
    \item for every $E \in \mathcal{T}$, there is a distinguished triangle
    \begin{displaymath}
        \tau^{\leq 0} E \to E \to \tau^{\geq 1} E \to (\tau^{\leq 0} E)[1]
    \end{displaymath}
    with $\tau^{\leq 0} E \in \mathcal{T}^{\leq 0}$ and $\tau^{\geq 1} E \in \mathcal{T}^{\geq 0}[-1]$.
\end{itemize}
This notion was introduced in \cite{Beilinson/Berstein/Deligne/Gabber:2018}. 
The above distinguished triangle is called the \textbf{truncation triangle} of $E$ with respect to the $t$-structure. 
For any $n \in \mathbb{Z}$, $(\mathcal{T}^{\leq n}, \mathcal{T}^{\geq n})$ is also a $t$-structure on $\mathcal{T}$ where $\mathcal{T}^{\leq n} := \mathcal{T}^{\leq 0}[-n]$ and $\mathcal{T}^{\geq n} := \mathcal{T}^{\geq 0}[-n]$. 
Define the subcategories:
\begin{displaymath}
    \mathcal{T}^{-}:= \bigcup_{n=0}^\infty \mathcal{T}^{\leq n},\quad \mathcal{T}^{+} := \bigcup^\infty_{n=0} \mathcal{T}^{\geq -n},\quad \mathcal{T}^b := \mathcal{T}^{-} \cap \mathcal{T}^{+}.
\end{displaymath}
Let $F \colon \mathcal{T} \to \mathcal{T}^\prime$ be an exact functor between triangulated categories equipped with $t$-structures $(\mathcal{T}^{\leq 0}, \mathcal{T}^{\geq 0})$ and $((\mathcal{T}^\prime)^{\leq 0}, (\mathcal{T}^\prime)^{\geq 0})$. 
We say that $F$ is \textbf{right $t$-exact} if $
F(\mathcal{T}^{\leq 0}) \subseteq (\mathcal{T}^\prime)^{\leq 0}$,
and \textbf{left $t$-exact} if 
$F(\mathcal{T}^{\geq 0}) \subseteq (\mathcal{T}^\prime)^{\geq 0}$. 
If both conditions hold, then $F$ is \textbf{$t$-exact}.

A strictly full subcategory $\mathcal{A} \subseteq \mathcal{T}$ is an \textbf{aisle} if it is closed under positive shifts and extensions, and the inclusion $\mathcal{A} \to \mathcal{T}$ admits a right adjoint. 
Given a $t$-structure $(\mathcal{T}^{\leq 0}, \mathcal{T}^{\geq 0})$, the subcategory $\mathcal{T}^{\leq 0}$ is an aisle. 
By \cite[\S 1, Proposition]{Keller/Vossieck:1988}, any aisle $\mathcal{A}$ determines a $t$-structure $(\mathcal{A}, \mathcal{A}^\perp[1])$, where
\begin{displaymath}
    \mathcal{A}^\perp := \{ T \in \mathcal{T} : \operatorname{Hom}(A,T) = 0 \text{ for all } A \in \mathcal{A} \}.
\end{displaymath}
The subcategories $\mathcal{T}^{\leq 0}$ and $\mathcal{T}^{\geq 0}$ are called the \textbf{aisle} and \textbf{coaisle} of the $t$-structure, respectively. 
More generally, a \textbf{preaisle} is a strictly full subcategory closed under positive shifts and extensions. 
By \cite[Lemma 1.3 \& Corollary 1.4]{AlonsoTarrio/Lopez/Salorio:2003}, aisles are closed under direct summands and under small coproducts when these exist in $\mathcal{T}$.

Assume $\mathcal{T}$ admits small coproducts. 
A preaisle $\mathcal{A} \subseteq \mathcal{T}$ is called \textbf{cocomplete} if it is closed under all coproducts in $\mathcal{T}$. 
Given $\mathcal{S}\subseteq \mathcal{T}$, $\overline{\langle \mathcal{S} \rangle}^{(-\infty, 0]}$ is defined to be the smallest cocomplete preaisle containing $\mathcal{S}$. 
By \cite[Theorem 2.3]{Neeman:2021b}, if $\mathcal{T}$ is well generated (e.g.\ compactly generated), $\overline{\langle \mathcal{S} \rangle}^{(-\infty, 0]}$ is an aisle whenever $\mathcal{S}$ is essentially small. 
Furthermore, if $\mathcal{A}\subseteq \mathcal{T}^c$ is closed under positive shifts, then \cite[Theorem 2.3.3 \& Remark 2.3.4]{Canonaco/Haesemeyer/Neeman/Stellari:2024} (which generalizes \cite[Theorem A.1 \& Proposition A.2]{AlonsoTarrio/Lopez/Salorio:2003}) says that $\operatorname{Coprod}(\mathcal{A})$ is the aisle $\overline{\langle \mathcal{A} \rangle}^{(-\infty, 0]}$. 
An aisle $\mathcal{U}$ on $\mathcal{T}$ is \textbf{compactly generated} when there exists a collection of compact objects $\mathcal{Q} \subseteq \mathcal{T}^c$ satisfying $\overline{\langle \mathcal{Q} \rangle}^{(-\infty, 0]} = \mathcal{U}$. 
In this case, we say $\mathcal{U}$ is compactly generated by $\mathcal{Q}$. 
We say a $t$-structure is \textbf{compactly generated} if its aisle is as such. 
For instance, if $G \in \mathcal{T}^c$, then 
\begin{displaymath}
    \operatorname{Coprod}( \{ G[n] : n\geq 0 \}) = \overline{\langle G \rangle}^{(-\infty, 0]} .
\end{displaymath}

Consider a pair of $t$-structures $(\mathcal{T}_1^{\leq 0},\mathcal{T}_1^{\geq 0})$ and $(\mathcal{T}_2^{\leq 0},\mathcal{T}_2^{\geq 0})$ on $\mathcal{T}$.
We say they are \textbf{equivalent} if there exists $A\geq 0$ such that $\mathcal{T}_1^{\leq -A}\subseteq \mathcal{T}_2^{\leq 0}\subseteq \mathcal{T}_1^{\leq A}$.
If $\mathcal{T}$ has a single compact generator, the \textbf{preferred equivalence class} is the equivalence class of the $t$-structure compactly generated by any compact generator.
By \cite[Remark 1.19]{Neeman:2026}, the preferred equivalence class does not depend on the choice of compact generator. 

If $\mathcal{T}$ has small coproducts and is equipped with a $t$-structure $(\mathcal{T}^{\leq 0}, \mathcal{T}^{\geq 0})$, let $\mathcal{T}^-_c$ be the collection of $E\in \mathcal{T}$ such that for every $n>0$ there exists a distinguished triangle 
\begin{displaymath}
    D \to E \to F \to D[1]
\end{displaymath}
with $D\in \mathcal{T}^c$ and $F\in \mathcal{T}^{\leq -n}$. 
Moreover, let $\mathcal{T}^b_c := \mathcal{T}^b \cap \mathcal{T}^-_c$. 
Furthermore, if $\mathcal{T}$ is weakly approximable and $(\mathcal{T}^{\leq 0}, \mathcal{T}^{\geq 0})$ is in the preferred equivalence class, then $\mathcal{T}^?$ and $\mathcal{T}^\#_c$ are independent of the $t$-structure representing the equivalence class for $?\in \{-,+,b\}$ and $\# \in \{-,b\}$. 
See \cite[Remark 4.1.5]{Canonaco/Haesemeyer/Neeman/Stellari:2024}.

%%%%%%%%%%%%%%%%%%%%%%%%%%%%%%%%%%%%%%%%%%%
\subsubsection{Recollements}
\label{sec:prelim_recollements}
%%%%%%%%%%%%%%%%%%%%%%%%%%%%%%%%%%%%%%%%%%%

We briefly recall the notion of a recollement. 
See \cite[\S 1.4]{Beilinson/Berstein/Deligne/Gabber:2018} for details. 
A \textbf{recollement} is a diagram of triangulated categories and exact functors of the form 
\begin{displaymath}
        \begin{tikzcd}[column sep=6em]
            {\mathcal{D}_Z} & {\mathcal{D}} & {\mathcal{D}_U}
            \arrow["{i_\ast}" description, from=1-1, to=1-2]
            \arrow["{i^\ast}"', bend right=35, from=1-2, to=1-1]
            \arrow["{i^!}", bend left=35, from=1-2, to=1-1]
            \arrow["{j^\ast}" description, from=1-2, to=1-3]
            \arrow["{j_!}"', bend right=35, from=1-3, to=1-2]
            \arrow["{j_\ast}", bend left=35, from=1-3, to=1-2]
        \end{tikzcd}
    \end{displaymath}
%% MISC: 'left' = \lambda and 'right' = \rho for adjoint functors!! Neeman suggested this idea
satisfying:
\begin{itemize}
    \item $i^\ast \dashv i_\ast \dashv i^!$ and $j_! \dashv j^\ast \dashv j_\ast$ (i.e.\ adjoint triples)
    \item $i_\ast, j_!, j_\ast$ are fully faithful
    \item $\ker (j^\ast)$ coincides with the strictly full subcategory on objects of the form $i_\ast(T)$ where $T\in \mathcal{D}_Z$.
\end{itemize}
In such a case, there are distinguished triangles
\begin{displaymath}
    \begin{aligned}
        j_! j^\ast E 
        & \xrightarrow{\epsilon_E} E 
        \xrightarrow{\eta_E} i_\ast i^\ast E 
        \xrightarrow{\varphi_E} j_! j^\ast E[1],
        \\& i_\ast i^! E \to E \to j_\ast j^\ast E \to i_\ast i^! E[1]
    \end{aligned}
\end{displaymath}
which are functorial in $E\in \mathcal{D}$.
In particular, the natural transformations between these functors are given by the (co)units of the relevant adjoint pairs. 
Since $j_!$, $j^\ast$, $i_\ast$, and $i^\ast$ are left adjoints, they preserve coproducts. 
Additionally, if each triangulated category is compactly generated, $i_\ast$ and $j^\ast$ admitting right adjoints implies $i^\ast$ and $j_!$ preserve compact objects \cite[Theorem 5.1]{Neeman:1996}.
%%MISC: One way to digest these facts, as suggested by Neeman, is that $\lambda$ starts with `L', which insinuates `left' adjoint.
Given $E\in \mathcal{D}_Z$ and $B \in \mathcal{D}_U$, $j^\ast i_\ast E \cong 0$, $i^! j_\ast B \cong 0$, and $i^\ast j_! B \cong 0$. 

Assume $(\mathcal{D}_Z^{\leq 0}, \mathcal{D}_Z^{\geq 0})$ and $(\mathcal{D}_U^{\leq 0}, \mathcal{D}_U^{\geq 0})$ are $t$-structures respectively on $\mathcal{D}_Z$ and $\mathcal{D}_U$. 
It is possible to `glue' these aisles to obtain a new $t$-structure on $\mathcal{D}$. 
Specifically, the glued $t$-structure is given by 
\begin{displaymath}
    \begin{aligned}
        \mathcal{D}^{\leq 0} &\colonequals  \left\{ E\in \mathcal{D} : i^\ast E \in \mathcal{D}_Z^{\leq 0}, j^\ast E \in \mathcal{D}_U^{\leq 0} \right\},
        \\ \mathcal{D}^{\geq 0} &\colonequals  \left\{ E\in \mathcal{D} : i^! E \in \mathcal{D}_Z^{\geq 0}, j^\ast E \in \mathcal{D}_U^{\geq 0} \right\}.
    \end{aligned}
\end{displaymath}
Moreover, a few general facts include: $i^\ast$ and $j_!$ are right $t$-exact;  $i_\ast$ and $j^\ast$ are $t$-exact; and $i^!$ and $j_\ast$ are left $t$-exact.
See \cite[1.4.10 \& 1.3.17.(iii)]{Beilinson/Berstein/Deligne/Gabber:2018} for details.

%%%%%%%%%%%%%%%%%%%%%%%%%%%%%%%%%%%%%%%%%%%%
\subsubsection{Weak approximability}
\label{sec:prelim_weak_approx}
%%%%%%%%%%%%%%%%%%%%%%%%%%%%%%%%%%%%%%%%%%%%  

We recall the definition of weak approximability for triangulated categories.
See \cite[Definition 1.25]{Neeman:2026}.

\begin{definition}
    \label{def:weak_approx}
    A triangulated category $\mathcal{T}$ is \textbf{weakly approximable} if there exist a compact generator $G$, a $t$-structure $(\mathcal{T}^{\leq 0},\mathcal{T}^{\geq 0})$ and an integer $A>0$ such that:
    \begin{enumerate}
        \item \label{def:weak_approx1} $G\in \mathcal{T}^{\leq A}$ and $\operatorname{Hom}(G,\mathcal{T}^{\leq -A})=0$;
        \item \label{def:weak_approx2} every $F\in \mathcal{T}^{\leq 0}$ admits a distinguished triangle $E\to F\to D\to E[1]$ with $E\in \overline{\langle G\rangle}^{[-A,A]}$ and $D\in \mathcal{T}^{\leq -1}$.
    \end{enumerate}
    We call $(G,\mathcal{T}^{\leq 0},A)$ \textbf{weak approximation data}.
\end{definition}

In our convention, compactly generated triangulated categories a priori have small coproducts.
In \Cref{def:weak_approx}, \eqref{def:weak_approx1} is the condition $G[A]\in \mathcal{T}^{\leq 0}$, $\operatorname{Hom}(G[-A],\mathcal{T}^{\leq 0})=0$ of \cite[2.9]{Burke/Neeman/Pauwels:2023}, since $\operatorname{Hom}(G[-A],\mathcal{T}^{\leq 0}) = \operatorname{Hom}(G,\mathcal{T}^{\leq -A})$.

A weakly approximable triangulated category $\mathcal{T}$ is \textbf{coherent} if, for any $t$-structure $(\mathcal{T}^{\leq 0}, \mathcal{T}^{\geq 0})$ in the preferred equivalence class, there exists an integer $N > 0$ such that every object $Y \in \mathcal{T}^-_c$ admits a distinguished triangle 
\begin{displaymath}
    X \to Y \to Z \to X[1]
\end{displaymath}
with $X \in \mathcal{T}^-_c \cap \mathcal{T}^{\leq N}$ and $Z\in \mathcal{T}^b_c \cap \mathcal{T}^{\geq 0}$.
See \cite[Definition 10.1.1]{Canonaco/Haesemeyer/Neeman/Stellari:2024}.

%%%%%%%%%%%%%%%%%%%%%%%%%%%%%%%%%%%%%
\subsection{Algebraic spaces}
\label{sec:algebraic_stacks}
%%%%%%%%%%%%%%%%%%%%%%%%%%%%%%%%%%%%%

We follow the conventions of \cite{StacksProject} and \cite[\S 2.2]{GuisadoVillaalgordo/Lank:2026} for algebraic spaces. 
Fix a quasi-compact quasi-separated algebraic space $X$. 
Write $|X|$ for its underlying topological space.

%%%%%%%%%%%%%%%%%%%%%%%%%%%%%%%%%%%%%%%%%%%%
\subsubsection{Categories/functors}
\label{sec:prelim_stacks_categories_functors}
%%%%%%%%%%%%%%%%%%%%%%%%%%%%%%%%%%%%%%%%%%%%

$\operatorname{Mod}(X)$ denotes the Grothendieck abelian category of sheaves of $\mathcal{O}_X$-modules on the small \'{e}tale site of $X$. 
$\operatorname{Qcoh}(X)$ is the strictly full subcategory (i.e.\ full and closed under isomorphisms) of $\operatorname{Mod}(X)$ consisting of quasi-coherent sheaves. 
$D(X) := D(\operatorname{Mod}(X))$ is the derived category of $\operatorname{Mod}(X)$.
$D_{\operatorname{qc}}(X)$ is the full subcategory of $D(X)$ consisting of complexes with quasi-coherent cohomology. 
$\operatorname{Perf}(X)$ denote the full subcategory of perfect complexes in $D_{\operatorname{qc}}(X)$. 
If $X$ is Noetherian, then $D^?_{\operatorname{coh}}(X)$ for complexes with coherent cohomology where $?\in \{-,+,b\}$. 
For a morphism $f\colon Y\to X$ of quasi-compact quasi-separated algebraic spaces, we use the adjunction $\mathbf{L}f^\ast\dashv\mathbf{R}f_\ast$ on $D_{\operatorname{qc}}$.

%%%%%%%%%%%%%%%%%%%%%%%%%%%%%%%%%%%%%%%%%%%%
\subsubsection{Support}
\label{sec:prelim_stacks_support}
%%%%%%%%%%%%%%%%%%%%%%%%%%%%%%%%%%%%%%%%%%%%

For $M \in \operatorname{Qcoh}(X)$, let $\operatorname{supp}(M) := p\big(\operatorname{supp}(p^\ast M)\big) \subseteq |X|$ where $p \colon U \to X$ is any \'{e}tale presentation. 
This definition is independent of the choice of $p$. 
More generally, given $E \in D_{\operatorname{qc}}(X)$, let
\begin{displaymath}
    \operatorname{supp}(E) := \bigcup_{j \in \mathbb{Z}} \operatorname{supp}\big(\mathcal{H}^j(E)\big) \subseteq |X|.
\end{displaymath}
This is called the \textbf{support} of $E$. 
Given an open immersion $j\colon U \to X$, $D_{\operatorname{qc},Z}(X):=\ker (\mathbf{L}j^\ast)$ where $Z=|X|\setminus |U|$. 
It can be checked that
\begin{displaymath}
    D_{\operatorname{qc},Z}(X) = \{ E \in D_{\operatorname{qc}}(X) : \operatorname{supp}(E)\subseteq Z \}.
\end{displaymath}
We use the analogous notation with coherent cohomology.

%%%%%%%%%%%%%%%%%%%%%%%%%%%%%%%%%%%%%%%%%%
\section{Auxiliary results}
\label{sec:auxillary}
%%%%%%%%%%%%%%%%%%%%%%%%%%%%%%%%%%%%%%%%%%%

\begin{lemma}
    \label{lem:closures}
    Let $\mathcal{T}$ and $\mathcal{T}^\prime$ be triangulated categories which admit small coproducts. 
    Choose $H\in \mathcal{T}$ and integers $a\leq b$.
    Then:
    \begin{enumerate}
        \item \label{lem:closures1} $\overline{\langle H\rangle}^{[a,b]}$ is closed under small coproducts, extensions, and direct summands; in fact, it is the smallest strictly full subcategory of $\mathcal{T}$ closed under these operations that contains $H[a,b]$.
        \item \label{lem:closures2} $\overline{\langle H\rangle}^{[a,b]}[s] = \overline{\langle H\rangle}^{[a-s,b-s]}$ and $\overline{\langle H\rangle}^{[a,b]} \subseteq \overline{\langle H\rangle}^{[a^\prime,b^\prime]}$ for all $s\in \mathbb{Z}$, $a^\prime\leq a$, and $b\leq b^\prime$.
        \item \label{lem:closures3} $F(\overline{\langle H\rangle}^{[a,b]}) \subseteq \overline{\langle H^\prime\rangle}^{[a+c,b+d]}$ for any small coproduct preserving exact $F\colon \mathcal{T}\to \mathcal{T}^\prime$ and $H^\prime \in \mathcal{T}^\prime$ with $F(H)\in \overline{\langle H^\prime\rangle}^{[c,d]}$.
        \item \label{lem:closures4} Let $Y\in \mathcal{T}$. If $\operatorname{Hom}(H[-k],Y)=0$ for all $a\leq k\leq b$, then $\operatorname{Hom}(X,Y)=0$ for all $X\in \overline{\langle H\rangle}^{[a,b]}$. If $H$ is compact, then $\{Y\in \mathcal{T} : \operatorname{Hom}(H,Y)=0\}$ is closed under small coproducts, extensions, and direct summands.
        \item \label{lem:closures5} $\operatorname{smd}(\operatorname{Coprod}(H[a,b])) \subseteq \operatorname{Coprod}(H[a,b+1])$.
    \end{enumerate}
\end{lemma}

\begin{proof}
    First we prove \eqref{lem:closures1}.
    By definition, $\overline{\langle H\rangle}^{[a,b]} := \operatorname{smd}(\operatorname{Coprod}(H[a,b]))$. 
    Let $\Lambda$ be a set.
    Choose $X_\lambda \in \overline{\langle H\rangle}^{[a,b]}$ where $\lambda \in \Lambda$. 
    For each $\lambda \in \Lambda$, there exists $X^\prime_\lambda \in \mathcal{T}$ such that $X_\lambda \oplus X^\prime_\lambda\in \operatorname{Coprod}(H[a,b])$.
    Since $\operatorname{Coprod}(H[a,b])$ is closed under small coproducts, $\oplus_{\lambda\in \Lambda} (X_\lambda \oplus X^\prime_\lambda)\in \operatorname{Coprod}(H[a,b])$. 
    Note that $\oplus_{\lambda\in \Lambda}X_\lambda$ is a direct summand of $\oplus_{\lambda\in \Lambda} (X_\lambda \oplus X^\prime_\lambda)$. 
    Hence, $\overline{\langle H\rangle}^{[a,b]}$ is closed under small coproducts. 

    Consider a distinguished triangle 
    \begin{displaymath}
        A \to B \to C \to A[1]
    \end{displaymath}
    where $A,C\in \overline{\langle H\rangle}^{[a,b]}$. 
    We can find $A^\prime,C^\prime$ such that $A\oplus A^\prime, C\oplus C^\prime \in \operatorname{Coprod}(H[a,b])$. 
    There exist distinguished triangles
    \begin{displaymath}
        A^\prime \xrightarrow{1_{A^\prime}} A^\prime \to 0 \to A^\prime[1]
    \end{displaymath}
    and \begin{displaymath}
        0 \to C^\prime \xrightarrow{1_{C^\prime}} C^\prime \to 0[1].
    \end{displaymath}
    Taking the direct sum of these distinguished triangles yields a distinguished triangle 
    \begin{displaymath}
        \begin{aligned}
            A\bigoplus A^\prime \bigoplus 0 
            \to B \bigoplus A^\prime \bigoplus C^\prime
            \to C \bigoplus 0 \bigoplus C^\prime \to A[1]\bigoplus A^\prime[1] \bigoplus 0.
        \end{aligned}
    \end{displaymath}
    See \cite[Proposition 1.2.1]{Neeman:2001}. 
    Since $A\oplus A^\prime, C\oplus C^\prime\in  \operatorname{Coprod}(H[a,b])$, it follows that $B \oplus A^\prime \oplus C^\prime\in \operatorname{Coprod}(H[a,b])$ because the category is closed under extensions.
    Consequently, $B \in \overline{\langle H\rangle}^{[a,b]}$.

    To see that $\overline{\langle H\rangle}^{[a,b]}$ is closed under direct summands, use the fact that a direct summand of a direct summand remains a direct summand.

    Lastly, we show the minimality claim. 
    Let $\mathcal{B}$ be a strictly full subcategory of $\mathcal{T}$ closed under small coproducts, extensions, and direct summands that contains $H[a,b]$.
    In this case, $\operatorname{Coprod}(H[a,b])\subseteq \mathcal{B}$. 
    Indeed, $\mathcal{B}$ is closed under small coproducts and extensions, and it contains $H[a,b]$. 
    Moreover, $\operatorname{smd}(\operatorname{Coprod}(H[a,b]))\subseteq \mathcal{B}$ because $\mathcal{B}$ is closed under direct summands. 

    Second we prove \eqref{lem:closures2}. 
    The autoequivalence $[s]\colon \mathcal{T} \to \mathcal{T}$ preserves extensions, small coproducts, and direct summands. 
    Hence, $\overline{\langle H\rangle}^{[a,b]}[s]$ is the smallest strictly full subcategory containing $(H[a,b])[s]$ which is closed under extensions, small coproducts, and direct summands.
    However, shifting by $s$ yields 
    \begin{displaymath}
        \begin{aligned}
            (H[a,b])[s] 
            &= \{ H [-k][s] : a\leq k \leq b \} 
            \\&= \{ H [-k+s] : a\leq k \leq b \}
            \\&=\{ H [-l] : a-s\leq l \leq b-s \}
            \\&=: H[a-s,b-s].
        \end{aligned}
    \end{displaymath}
    It follows that $\overline{\langle H\rangle}^{[a,b]}[s] = \overline{\langle H\rangle}^{[a-s,b-s]}$. 
    Moreover, as $H[a,b]\subseteq H[a^\prime,b^\prime]$, \eqref{lem:closures1} implies $\overline{\langle H\rangle}^{[a,b]} \subseteq \overline{\langle H\rangle}^{[a^\prime,b^\prime]}$.

    Third we prove \eqref{lem:closures3}.
    Choose $a\leq k\leq b$.
    Note that $F$ is exact and $F(H)\in \overline{\langle H^\prime\rangle}^{[c,d]}$.
    Hence, by \eqref{lem:closures2},
    \begin{displaymath}
        F(H[-k])\cong F(H)[-k]
        \in \overline{\langle H^\prime\rangle}^{[c,d]}[-k] 
        = \overline{\langle H^\prime\rangle}^{[c+k,d+k]}\subseteq \overline{\langle H^\prime\rangle}^{[a+c,b+d]}.
    \end{displaymath}
    Set $\mathcal{S}$ to be the full subcategory of $E\in \mathcal{T}$ such that $F(E)\in \overline{\langle H^\prime\rangle}^{[a+c,b+d]}$.
    It follows that $H[a,b]\subseteq \mathcal{S}$.
    By \eqref{lem:closures1}, $\overline{\langle H^\prime\rangle}^{[a+c,b+d]}$ is strictly full and closed under small coproducts, extensions, and direct summands.
    Since $F$ preserves small coproducts, $\mathcal{S}$ has the same properties.
    By minimality, \eqref{lem:closures1} implies $\overline{\langle H\rangle}^{[a,b]}\subseteq \mathcal{S}$.
    Thus, \eqref{lem:closures3} follows.

    Next we prove \eqref{lem:closures4}. 
    It is easy to see that $\{E : \operatorname{Hom}(E,Y)=0\}$ is strictly full. 
    Let $\Lambda$ be a set. 
    Choose $X_\lambda \in \{E : \operatorname{Hom}(E,Y)=0\}$.
    By the universal property of coproducts, there exists an isomorphism
    \begin{displaymath}
        \operatorname{Hom}(\bigoplus_{\lambda \in \Lambda} X_\lambda, Y) 
        \cong \prod_{\lambda \in \Lambda} \operatorname{Hom}(X_\lambda, Y)
        \cong 0.
    \end{displaymath} 
    See \cite[Example 1.1.13]{Neeman:2001} applied to $\mathcal{T}^{op}$. 
    The last isomorphism follows from the fact that each $X_\lambda \in \{E : \operatorname{Hom}(E,Y)=0\}$.
    Hence, $\{E : \operatorname{Hom}(E,Y)=0\}$ is closed under small coproducts. 
    Let $A,B\in \{E : \operatorname{Hom}(E,Y)=0\}$.
    Consider an extension $C$ given by a distinguished triangle 
    \begin{displaymath}
        A \to C \to B \to A[1].
    \end{displaymath}
    Apply $\operatorname{Hom}(-,Y)$ to obtain a long exact sequence in Hom sets 
    \begin{displaymath}
        \cdots \to \operatorname{Hom}(B,Y) \to \operatorname{Hom}(C,Y) \to \operatorname{Hom}(A,Y) \to \cdots.
    \end{displaymath}
    Hence, $\operatorname{Hom}(C,Y) \cong 0$ if 
    \begin{displaymath}
        \operatorname{Hom}(A,Y) \cong \operatorname{Hom}(B,Y) \cong 0.
    \end{displaymath}
    This implies $\{E : \operatorname{Hom}(E,Y)=0\}$ is closed under extensions.
    Moreover, if $A\oplus B \in \{E : \operatorname{Hom}(E,Y)=0\}$, then 
    \begin{displaymath}
        \operatorname{Hom}(A\oplus B,Y)\cong \operatorname{Hom}(A,Y) \bigoplus \operatorname{Hom}(B,Y) \cong 0.
    \end{displaymath} 
    Thus, $\{E : \operatorname{Hom}(E,Y)=0\}$ is closed under direct summands. 
    The hypothesis asserts that $H[a,b]\subseteq \{E : \operatorname{Hom}(E,Y)=0\}$.
    By \eqref{lem:closures1}, $\overline{\langle H \rangle}^{[a,b]}\subseteq \{E : \operatorname{Hom}(E,Y)=0\}$.
    A similar argument with $H$ compact proves the second claim for \eqref{lem:closures4}.
    Indeed, the compactness of $H$ gives the closure under coproducts, whereas the other conditions are analogous.
    
    Lastly we prove \eqref{lem:closures5}. 
    Let $X\in \operatorname{smd}(\operatorname{Coprod}(H[a,b]))$.
    Choose $X^\prime$ such that $X\oplus X^\prime \in \operatorname{Coprod}(H[a,b])$.
    The so-called Eilenberg swindle yields an isomorphism
    \begin{displaymath}
        \begin{aligned}
            X \oplus (\bigoplus_{n\in \mathbb{N}} (X\oplus X^\prime)) 
            &\cong X \oplus (\bigoplus_{n\in \mathbb{N}} (X^\prime \oplus X)) 
            \\&\cong X \oplus (X^\prime \oplus X) \oplus (X^\prime \oplus X) \oplus \cdots
            \\&\cong (X \oplus X^\prime) \oplus (X \oplus X^\prime) \oplus \cdots
            \\&\cong \bigoplus_{n\in \mathbb{N}} (X \oplus X^\prime).
        \end{aligned}
    \end{displaymath}
    Consequently, there exists a distinguished triangle 
    \begin{displaymath}
        X \to X \oplus (\bigoplus_{n\in \mathbb{N}} (X\oplus X^\prime)) \to \bigoplus_{n\in \mathbb{N}}(X\oplus X^\prime) \to X[1].
    \end{displaymath}
    %%NOTE: A\to A\oplus B \to B \to A[1].
    Rotating backwards yields 
    \begin{displaymath}
        (\bigoplus_{n\in \mathbb{N}}( X\oplus X^\prime))[-1] 
        \to X \to X \oplus (\bigoplus_{n\in \mathbb{N}} (X\oplus X^\prime)) \to \bigoplus_{n\in \mathbb{N}}( X\oplus X^\prime).
    \end{displaymath}
    Since $[a,b]\subseteq [a,b+1]$, $\operatorname{Coprod}(H[a,b]) \subseteq \operatorname{Coprod}(H[a,b+1])$.
    Observe that 
    \begin{displaymath}
        \begin{aligned}
            (\bigoplus_{n\in \mathbb{N}} (X\oplus X^\prime))[-1]
            &\in \operatorname{Coprod}(H[a,b])[-1] 
            \\&= \operatorname{Coprod}(H[a+1,b+ 1])
            \\&\subseteq \operatorname{Coprod}(H[a,b+1]).
        \end{aligned}
    \end{displaymath}
    Therefore, 
    \begin{displaymath}
        \begin{aligned}
            X 
            &\in \operatorname{Coprod}(H[a,b+1]) \star \operatorname{Coprod}(H[a,b])
            \\&\subseteq \operatorname{Coprod}(H[a,b+1]) \star \operatorname{Coprod}(H[a,b+1])
            \\&\subseteq \operatorname{Coprod}(H[a,b+1]).
        \end{aligned}
    \end{displaymath}
    This completes the proof.
\end{proof}

\begin{remark}
    \label{rmk:octahedral}
    We use the octahedral axiom in the following ways.
    Let $X\xrightarrow{f} Y\xrightarrow{g} Z$ be composable.
    \begin{enumerate}
        \item \label{rmk:octahedraloc} Choose distinguished triangles $X\xrightarrow{f} Y\to C_f\to X[1]$, $Y\xrightarrow{g} Z\xrightarrow{p_g} C_g\to Y[1]$ and $X\xrightarrow{gf} Z\xrightarrow{p_{gf}} C_{gf}\to X[1]$. There is a distinguished triangle $C_f\to C_{gf}\xrightarrow{h} C_g\to C_f[1]$ with $h\circ p_{gf} = p_g$. See e.g.\ \cite[Proposition 1.4.6]{Neeman:2001}.
        \item \label{rmk:octahedralof} Choose distinguished triangles $F_f\to X\xrightarrow{f} Y\to F_f[1]$, $F_g\to Y\xrightarrow{g} Z\to F_g[1]$ and $F_{gf}\to X\xrightarrow{gf} Z\to F_{gf}[1]$. There is a distinguished triangle $F_f\to F_{gf}\to F_g\to F_f[1]$. Indeed, take the distinguished triangle $C_f\to C_{gf}\xrightarrow{h} C_g\to C_f[1]$ and shift by $[-1]$. Here we use that $F_? \cong C_?[-1]$.
    \end{enumerate}
\end{remark}

\begin{lemma}
    \label{lem:aisle}
    Let $(\mathcal{T}^{\leq 0},\mathcal{T}^{\geq 0})$ be a $t$-structure on a triangulated category $\mathcal{T}$.
    Consider a distinguished triangle  
    \begin{displaymath}
        E\to F\to D\to E[1]
    \end{displaymath}
    where $F\in \mathcal{T}^{\leq 0}$ and $D\in \mathcal{T}^{\leq -m}$ for some $m\geq 1$.
    Then $E\in \mathcal{T}^{\leq 0}$.
\end{lemma}

\begin{proof}
    By rotating backwards, there exists a distinguished triangle  
    \begin{displaymath}
       D[-1] \to E\to F\to D.
    \end{displaymath}
    Note that $D[-1]\in \mathcal{T}^{\leq 1-m} \subseteq \mathcal{T}^{\leq 0}$.
    Since aisles are closed under extensions, $E\in \mathcal{T}^{\leq 0}$.
\end{proof}

\begin{lemma}
    \label{lem:iterate}
    Let $(H,\mathcal{T}^{\leq 0},A)$ be weak approximation data for $\mathcal{T}$.
    For every $F\in \mathcal{T}^{\leq 0}$ and every $m\geq 1$, there exists a distinguished triangle 
    \begin{displaymath}
        E_m\to F\to D_m\to E_m[1]
    \end{displaymath}
    with $E_m\in \overline{\langle H\rangle}^{[1-m-A,A]} \cap \mathcal{T}^{\leq 0}$ and $D_m\in \mathcal{T}^{\leq -m}$. 
\end{lemma}

\begin{proof}
    This follows from \cite[Corollary 3.2(1)]{Neeman:2026} with weak approximation data for $\mathcal{T}$.
    To be self-contained, we give a proof.
    By \Cref{lem:aisle}, for any desired distinguished triangle 
    \begin{displaymath}
        E_m\to F\to D_m\to E_m[1],
    \end{displaymath}
    it follows that $E_m\in \mathcal{T}^{\leq 0}$.
    Hence, it suffices to produce the desired distinguished triangle.
    We argue by induction on $m$. 
    The base case $m=1$ is \Cref{def:weak_approx}\eqref{def:weak_approx2}.

    Assume that the desired distinguished triangle exists for $m$, i.e.\ a distinguished triangle 
    \begin{displaymath}
        E_m\to F\xrightarrow{f} D_m\to E_m[1]
    \end{displaymath}
    with $E_m\in \overline{\langle H\rangle}^{[1-m-A,A]} \cap \mathcal{T}^{\leq 0}$ and $D_m\in \mathcal{T}^{\leq -m}$. 
    Since $D_m[-m]\in \mathcal{T}^{\leq 0}$, \Cref{def:weak_approx}\eqref{def:weak_approx2} yields a distinguished triangle 
    \begin{displaymath}
        E^\prime \to D_m[-m]\to D^\prime \to E^\prime[1]
    \end{displaymath}
    where $E^\prime \in \overline{\langle H\rangle}^{[-A,A]}$ and $D^\prime \in \mathcal{T}^{\leq -1}$.
    %%NOTE: Note that \Cref{lem:closures}\eqref{lem:closures1} ensures that $E^\prime \in \overline{\langle H\rangle}^{[-A,A]}$. 
    After shifting by $[m]$, we have a distinguished triangle 
    \begin{displaymath}
        E^\prime[m]\to D_m\xrightarrow{g} D^\prime[m]\to E^\prime[m+1]
    \end{displaymath}
    where $E^\prime[m]\in \overline{\langle H\rangle}^{[-A-m,A-m]}$ and $D^\prime[m]\in \mathcal{T}^{\leq -1-m}$.
    Set $D_{m+1} := D^\prime[m]$. 
    Complete the composite $g\circ f \colon F\to D^\prime[m]$ to a distinguished triangle,
    \begin{displaymath}
        F\to D^\prime[m] \to R \to F[1].
    \end{displaymath}
    Set $E_{m+1}:=R[-1]$. 
    The cocones of $f$, $g$, and $g\circ f$ are respectively $E_m$, $E^\prime[m]$, and $E_{m+1}$. 
    By \Cref{rmk:octahedral}\eqref{rmk:octahedralof}, this yields a distinguished triangle 
    \begin{displaymath}
        E_m \to E_{m+1} \to E^\prime [m] \to E_m [1].
    \end{displaymath}
    By \Cref{lem:closures}\eqref{lem:closures1}\eqref{lem:closures2},
    \begin{displaymath}
        \begin{aligned}
            E_{m+1} 
            & \in \overline{\langle H\rangle}^{[1-m-A,A]} \star \overline{\langle H\rangle}^{[-m-A,A-m]}
            \\& \subseteq \overline{\langle  H\rangle}^{[-m-A,A]} 
            \\&= \overline{\langle H\rangle}^{[1-(m+1)-A,A]}.
        \end{aligned}
    \end{displaymath}
    This completes the proof.
\end{proof}

\begin{lemma}
    \label{lem:hom_vanishing}
    Let $\mathcal{T}$ be a triangulated category admitting small coproducts. 
    Choose $H\in \mathcal{T}^c$.
    Denote by $\mathcal{T}^{\leq 0}$ the aisle compactly generated by $H$.
    Assume that $\operatorname{Hom}(H,H[n])=0$ for all $n\geq n_0$.
    Then $\operatorname{Hom}(H,\mathcal{T}^{\leq -n_0})=0$.
\end{lemma}

\begin{proof}
    This is \cite[Remark 1.24]{Neeman:2026}; cf.\ \cite[Lemma 3.5(1)]{Burke/Neeman/Pauwels:2023}.
    On one hand, 
    \begin{displaymath}
        \begin{aligned}
            \mathcal{T}^{\leq -n_0} 
            &= \mathcal{T}^{\leq 0}[n_0] 
            \\&= \operatorname{Coprod}(\{H[n] : n\geq n_0\}).
        \end{aligned}
    \end{displaymath}
    On the other hand, the strictly full subcategory 
    \begin{displaymath}
        \{Y \in \mathcal{T} : \operatorname{Hom}(H,Y)=0\}
    \end{displaymath}
    contains $H[n]$ for $n\geq n_0$. 
    Moreover, by \Cref{lem:closures}\eqref{lem:closures4}, this subcategory above is closed under small coproducts and extensions.
    Consequently, the subcategory above contains $\mathcal{T}^{\leq -n_0}$.
    Thus, the desired claim follows.
\end{proof}

%%%%%%%%%%%%%%%%%%%%%%%%%%%%%%%%%%%%%%%%%%%
\section{Lifts to aisles}
\label{sec:lifts}
%%%%%%%%%%%%%%%%%%%%%%%%%%%%%%%%%%%%%%%%%%%

\begin{setup}
    \label{setup:recollement}
    Consider a recollement of singly compactly generated triangulated categories: 
    \begin{displaymath}
        \begin{tikzcd}[column sep=6em]
            {\mathcal{D}_Z} & {\mathcal{D}} & {\mathcal{D}_U.}
            \arrow["{i_\ast}" description, from=1-1, to=1-2]
            \arrow["{i^\ast}"', bend right=35, from=1-2, to=1-1]
            \arrow["{i^!}", bend left=35, from=1-2, to=1-1]
            \arrow["{j^\ast}" description, from=1-2, to=1-3]
            \arrow["{j_!}"', bend right=35, from=1-3, to=1-2]
            \arrow["{j_\ast}", bend left=35, from=1-3, to=1-2]
        \end{tikzcd}
    \end{displaymath}
    For every $E\in \mathcal{D}$, there are distinguished triangles
    \begin{equation}
        \label{eq:gluing_triangles}
        \begin{aligned}
            j_! j^\ast E 
            & \xrightarrow{\epsilon_E} E 
            \xrightarrow{\eta_E} i_\ast i^\ast E 
            \xrightarrow{\varphi_E} j_! j^\ast E[1],
            \\& i_\ast i^! E \to E \to j_\ast j^\ast E \to i_\ast i^! E[1].
        \end{aligned}
    \end{equation}
    Here $\epsilon$ is the counit of $j_! \dashv j^\ast$ and $\eta$ is the unit of $i^\ast \dashv i_\ast$.
\end{setup}

\begin{lemma}
    \label{lem:glued}
    Assume \Cref{setup:recollement}.
    Consider $t$-structures $(\mathcal{D}_Z^{\leq 0},\mathcal{D}_Z^{\geq 0})$ and $(\mathcal{D}_U^{\leq 0},\mathcal{D}_U^{\geq 0})$ respectively on $\mathcal{D}_Z$ and $\mathcal{D}_U$. 
    Denote by $(\mathcal{D}^{\leq 0},\mathcal{D}^{\geq 0})$ the glued $t$-structure.
    For every $n\in \mathbb{Z}$,
    \begin{displaymath}
        \mathcal{D}^{\leq n} = \{E\in \mathcal{D} :\ i^\ast E\in \mathcal{D}_Z^{\leq n}\ \text{and}\ j^\ast E\in \mathcal{D}_U^{\leq n}\}.
    \end{displaymath}
    In particular,
    \begin{displaymath}
        i^\ast \mathcal{D}^{\leq n}\subseteq \mathcal{D}_Z^{\leq n},
        \qquad 
        j^\ast \mathcal{D}^{\leq n}\subseteq \mathcal{D}_U^{\leq n},
        \qquad 
        i_\ast \mathcal{D}_Z^{\leq n}\subseteq \mathcal{D}^{\leq n},
        \qquad 
        j_! \mathcal{D}_U^{\leq n}\subseteq \mathcal{D}^{\leq n}.
    \end{displaymath}
\end{lemma}

\begin{proof}
    If $n=0$, then this is the definition of the glued aisle.
    Since $i^\ast$ and $j^\ast$ commute with shifts, 
    \begin{displaymath}
        \mathcal{D}^{\leq n} = \{E\in \mathcal{D} :\ i^\ast E\in \mathcal{D}_Z^{\leq n}\ \text{and}\ j^\ast E\in \mathcal{D}_U^{\leq n}\}.
    \end{displaymath}
    Moreover, the functors $i^\ast, j^\ast, i_\ast, j_!$ are right $t$-exact for the $n=0$ case.
    After shifting, the desired containments follow.
\end{proof}

\begin{lemma}
    \label{lem:normalize}
    Assume \Cref{setup:recollement} where $\mathcal{D}_U$ is weakly approximable with compact generator $G_U$.
    Suppose $G^\prime$ is a compact generator for $\mathcal{D}$ which satisfies $\operatorname{Hom}(G^\prime,G^\prime[n])=0$ for $n\gg 0$.
    Denote by $\mathcal{D}_U^{\leq 0}$ the aisle compactly generated by $G_U$.
    Then the following hold:
    \begin{enumerate}
        \item \label{lem:normalize1} there exists an $m\geq 0$ such that $j^\ast G^\prime \in \mathcal{D}_U^{\leq m}$
        \item \label{lem:normalize2} $G := G^\prime[m]\oplus j_! G_U$ is a compact generator of $\mathcal{D}$, and $i^\ast G$ is a compact generator of $\mathcal{D}_Z$
        \item \label{lem:normalize3} the $t$-structure on $\mathcal{D}$ compactly generated by $G$ is the gluing of the $t$-structure on $\mathcal{D}_Z$ compactly generated by $i^\ast G$ with $(\mathcal{D}_U^{\leq 0},\mathcal{D}_U^{\geq 0})$
        \item \label{lem:normalize4} $\operatorname{Hom}(G,G[n])=0$ for $n\gg 0$.
    \end{enumerate}
\end{lemma}

\begin{proof}
    \eqref{lem:normalize1} and \eqref{lem:normalize3} are essentially from \cite[Lemma 3.11]{Burke/Neeman/Pauwels:2023}.
    The observation below is \cite[Lemma 3.5(2)]{Burke/Neeman/Pauwels:2023} (see also \cite[Lemma A.1]{GuisadoVillaalgordo/Lank:2026}). 
    We recall the argument for convenience.

    We prove \eqref{lem:normalize1}.
    Recall that $j_!$ preserves compact objects.
    Hence, $j_! G_U\in \mathcal{D}^c$.
    We observe that for all $F,F^\prime\in \mathcal{D}^c$, there exists an integer $N$ such that $\operatorname{Hom}(F,F^\prime[n])=0$ for all $n\geq N$.
    First, consider the full subcategory $\mathcal{H}$ of $F\in \mathcal{D}^c$ where such an $N$ exists with $F^\prime = G^\prime$.
    This subcategory is thick (e.g.\ use the long exact sequences of $\operatorname{Hom}(-,G^\prime[n])$). 
    By hypothesis, $G^\prime \in \mathcal{H}$.
    As $G^\prime$ is a classical generator of $\mathcal{D}^c$ (see \Cref{sec:prelim_generation}), $\mathcal{H} = \mathcal{D}^c$.
    Second, fix $F\in \mathcal{D}^c$. 
    Consider the full subcategory $\mathcal{H}^\prime$ of $F^\prime\in \mathcal{D}^c$ for which such an $N$ exists.
    Again, $\mathcal{H}^\prime$ is thick, which can be shown using the long exact sequences of $\operatorname{Hom}(F,-)$.
    Also, by the first step above, $G^\prime \in \mathcal{H}^\prime$. 
    Hence, $\mathcal{H}^\prime = \mathcal{D}^c$, which proves the observation. 
    %%NOTE: cf.\ \cite[Lemma 3.5(2)]{Burke/Neeman/Pauwels:2023}.
    Applying the observation to $F = j_! G_U$ and $F^\prime = G^\prime$, and using adjunction, we have for $n\gg 0$,
    \begin{displaymath}
        \operatorname{Hom}(G_U[-n],j^\ast G^\prime)\cong \operatorname{Hom}(j_! G_U[-n],G^\prime) = 0.
    \end{displaymath}
    Since $\mathcal{D}_U$ is weakly approximable and $\mathcal{D}_U^{\leq 0}$ is in the preferred equivalence class, \cite[Corollary 3.10(1)]{Burke/Neeman/Pauwels:2023} implies that $j^\ast G^\prime\in \mathcal{D}_U^{\leq m}$ for some $m\geq 0$.

    Note that $j_! G_U$ is compact. 
    Hence, the hypotheses of \cite[Lemma A.1]{GuisadoVillaalgordo/Lank:2026} hold: $G^\prime$ and $G_U$ are respectively compact generators of $\mathcal{D}$ and $\mathcal{D}_U$, $\mathcal{D}_U^{\leq 0}$ is the aisle compactly generated by $G_U$, and $j^\ast G^\prime\in \mathcal{D}_U^{\leq m}$ by \eqref{lem:normalize1}.
    Consequently, \eqref{lem:normalize2} and \eqref{lem:normalize3} follow from \cite[Lemma A.1]{GuisadoVillaalgordo/Lank:2026}.

    To prove \eqref{lem:normalize4}, apply the observation with $F=F^\prime=G$.
\end{proof}

\begin{setup}
    \label{setup:standing}
    Consider \Cref{setup:recollement}.
    Assume that $\mathcal{D}_Z$ and $\mathcal{D}_U$ are weakly approximable.
    Suppose $\mathcal{D}$ has a compact generator $G^\prime$ satisfying $\operatorname{Hom}(G^\prime,G^\prime[n])=0$ for $n\gg 0$.
    Fix a compact generator $G_U$ of $\mathcal{D}_U$.
    Choose $m$ and $G = G^\prime[m]\oplus j_! G_U$ from \Cref{lem:normalize}.
    Consider the following $t$-structures:
    \begin{itemize}
        \item $\mathcal{D}_U^{\leq 0}$ is compactly generated by $G_U$;
        \item $\mathcal{D}_Z^{\leq 0}$ is compactly generated by $i^\ast G$;
        \item $\mathcal{D}^{\leq 0}$ is the glued $t$-structure (by \Cref{lem:normalize}\eqref{lem:normalize3}, it is the $t$-structure compactly generated by $G$; hence, $G\in \mathcal{D}^{\leq 0}$).
    \end{itemize}
    Note that $\mathcal{D}_U^{\leq 0}$ and $\mathcal{D}_Z^{\leq 0}$ are respectively in the preferred equivalence classes of $\mathcal{D}_U$ and $\mathcal{D}_Z$. 
    By \cite[Lemma 3.7]{Burke/Neeman/Pauwels:2023} (equivalently, \cite[Proposition 3.6]{Neeman:2026}), \Cref{lem:hom_vanishing}, and \Cref{lem:normalize}\eqref{lem:normalize4}, there exist $A_1,A_2,A_3 \gg 0$ such that the following hold:
    \begin{enumerate}
        \item \label{setup:standinga1} $(i^\ast G,\mathcal{D}_Z^{\leq 0},A_1)$ is weak approximation data for $\mathcal{D}_Z$
        \item \label{setup:standinga2} $(G_U,\mathcal{D}_U^{\leq 0},A_2)$ is weak approximation data for $\mathcal{D}_U$
        \item \label{setup:standinga3} $\operatorname{Hom}(G,\mathcal{D}^{\leq -A_3})=0$.
    \end{enumerate}
    Set $A= \max\{A_1,A_2,A_3\}$. 
    Define $K := A+1$, $B := 3A$, and $\mathcal{P} := \overline{\langle G\rangle}^{[-B,B]} \cap \mathcal{D}^{\leq K}$.
    By \Cref{lem:closures}\eqref{lem:closures1}, and the fact that aisles are closed under extensions and small coproducts, $\mathcal{P}$ is closed under small coproducts and extensions.
    Since $j_! G_U$ is a direct summand of $G$, it follows that $j_! G_U\in \overline{\langle G\rangle}^{[0,0]}$.
\end{setup}

\begin{definition}
    \label{def:lift}
    Assume \Cref{setup:standing}. 
    Let $E^\prime \in \mathcal{D}_Z$, $Y\in \mathcal{D}$, and $f^\prime \colon E^\prime \to i^\ast Y$.
    A \textbf{$\mathcal{P}$-lift} of $f^\prime$ is a pair $(q,\theta)$ consisting of a morphism $q\colon E\to Y$ with $E\in \mathcal{P}$ and an isomorphism $\theta\colon i^\ast E \to E^\prime$ such that $i^\ast q = f^\prime \circ \theta$.
\end{definition}

\begin{lemma}
    \label{lem:onecell}
    Assume \Cref{setup:standing}. 
    Let $Y\in \mathcal{D}^{\leq K}$, $-A-1\leq n\leq A$, and $f^\prime \colon (i^\ast G)[n]\to i^\ast Y$.
    Then $f^\prime$ admits a $\mathcal{P}$-lift.
\end{lemma}

\begin{proof}
    The following proof is a variant of \cite[Lemma 5.3]{Burke/Neeman/Pauwels:2023}.
    Specifically, the lift is constructed directly from the adjunct $g$ rather than from a roof \cite[Lemma 5.2]{Burke/Neeman/Pauwels:2023}, and it lies in $\mathcal{D}^{\leq K}$ for the same $K$ that bounds $Y$.

    Since $i^\ast$ is exact, $i^\ast(G[n])$ can be identified with $(i^\ast G)[n]$.
    Hence, $f^\prime$ yields a morphism $i^\ast(G[n])\to i^\ast Y$.
    Denote by $\eta$ the unit and $\overline{\varepsilon}$ the counit of $i^\ast\dashv i_\ast$.

    Set $M := j^\ast Y$.
    By \Cref{lem:glued}, $M\in \mathcal{D}_U^{\leq K}$, and so $M[K]\in \mathcal{D}_U^{\leq 0}$.
    Apply \Cref{lem:iterate} to the weak approximation data $(G_U,\mathcal{D}_U^{\leq 0},A)$ of \Cref{setup:standing}\eqref{setup:standinga2}, $M[K]$, and $3A+1 = K+2A$ in place of $m$.
    This gives a distinguished triangle
    \begin{displaymath}
        W^\prime\to M[K]\to V^\prime\to W^\prime[1].
    \end{displaymath}
    Here $W^\prime\in \overline{\langle G_U\rangle}^{[-4A,A]}\cap \mathcal{D}_U^{\leq 0}$ and $V^\prime\in \mathcal{D}_U^{\leq -3A-1}$.
    Define $W := W^\prime[-K]$ and $V := V^\prime[-K]$.
    Shifting by $[-K]$ yields a distinguished triangle
    \begin{equation}
        \label{eq:onecell_wc}
        W\xrightarrow{w} M\xrightarrow{c} V\to W[1].
    \end{equation}
    By \Cref{lem:closures}\eqref{lem:closures2},
    \begin{displaymath}
        W\in \overline{\langle G_U\rangle}^{[-4A+K,A+K]} = \overline{\langle G_U\rangle}^{[1-3A,2A+1]}.
    \end{displaymath}
    Moreover, $W\in \mathcal{D}_U^{\leq 0}[-K] = \mathcal{D}_U^{\leq K}$ and $V\in \mathcal{D}_U^{\leq -3A-1}[-K] = \mathcal{D}_U^{\leq -2A}$.

    Now let $g\colon G[n]\to i_\ast i^\ast Y$ be the adjunct of $f^\prime$.
    That is, $g = i_\ast(f^\prime)\circ \eta_{G[n]}$, and so $f^\prime = \overline{\varepsilon}_{i^\ast Y}\circ i^\ast(g)$.
    Set
    \begin{displaymath}
        \varphi := \varphi_Y\colon i_\ast i^\ast Y\to j_! M[1]
    \end{displaymath}
    to be the morphism in \eqref{eq:gluing_triangles}, and $\xi := \varphi\circ g$. 
    By \Cref{lem:glued}, $j_! V\in \mathcal{D}^{\leq -2A}$.
    As $n\leq A$, 
    \begin{displaymath}
        j_! V[1-n]\in \mathcal{D}^{\leq n-1-2A}\subseteq \mathcal{D}^{\leq -A}.
    \end{displaymath}
    Hence, \Cref{setup:standing}\eqref{setup:standinga3} gives
    \begin{equation}
        \label{eq:onecell_vanishing}
        \operatorname{Hom}(G[n],j_! V[1])\cong \operatorname{Hom}(G,j_! V[1-n]) = 0.
    \end{equation}
    Applying $j_!$ to \eqref{eq:onecell_wc} gives a distinguished triangle
    \begin{displaymath}
        j_! W\xrightarrow{j_!(w)} j_! M \xrightarrow{j_!(c)} j_! V\to j_! W[1].
    \end{displaymath}
    By \cite[Lemma 1.1.10]{Neeman:2001}, $\operatorname{Hom}(G[n-1],-)$ is a homological functor on $\mathcal{D}$.
    This yields an exact sequence
    \begin{displaymath}
        \operatorname{Hom}(G[n-1],j_! W)
        \xrightarrow{j_!(w)\circ -} \operatorname{Hom}(G[n-1],j_! M)
        \xrightarrow{j_!(c)\circ -} \operatorname{Hom}(G[n-1],j_! V).
    \end{displaymath}
    By \eqref{eq:onecell_vanishing}, 
    \begin{displaymath}
        \operatorname{Hom}(G[n-1],j_! V)\cong \operatorname{Hom}(G[n],j_! V[1]) = 0.
    \end{displaymath}
    Hence, the left-hand map is surjective.
    Thus, there exists $\xi^{\prime\prime}\colon G[n-1]\to j_! W$ such that $\xi[-1] = j_!(w)\circ \xi^{\prime\prime}$.
    Setting $\xi^\prime := \xi^{\prime\prime}[1]\colon G[n]\to j_! W[1]$, we get $\xi = j_!(w)[1]\circ \xi^\prime$.

    Next we construct the lift.
    Completing $\xi^\prime$ and rotating backwards twice yields a distinguished triangle
    \begin{equation}
        \label{eq:onecell_triangle}
        j_! W\xrightarrow{\alpha} Q\xrightarrow{\beta} G[n]\xrightarrow{\xi^\prime} j_! W[1].
    \end{equation}
    %%NOTE: exists by TR1 and TR2.
    Consider the diagram
    \begin{displaymath}
        \begin{tikzcd}
            j_! W \ar[r, "\alpha"] \ar[d, "j_!(w)"'] & Q \ar[r, "\beta"] \ar[d, dashed, "q"] & G[n] \ar[r, "\xi^\prime"] \ar[d, "g"] & j_! W[1] \ar[d, "{j_!(w)[1]}"] \\
            j_! M \ar[r, "\epsilon_Y"] & Y \ar[r, "\eta_Y"] & i_\ast i^\ast Y \ar[r, "\varphi"] & j_! M [1]
        \end{tikzcd}
    \end{displaymath}
    with distinguished rows.
    Indeed, the top row is \eqref{eq:onecell_triangle}, whereas the bottom row is the first triangle of \eqref{eq:gluing_triangles} for $E = Y$ (i.e.\ use that $M = j^\ast Y$).
    The right-hand square commutes because 
    \begin{displaymath}
        \varphi\circ g = \xi = j_!(w)[1]\circ \xi^\prime.
    \end{displaymath}
    By TR3, applied after rotating both rows backwards, there exists $q\colon Q\to Y$ making the diagram commute.
    In particular, $\eta_Y\circ q = g\circ \beta$.

    Set $\theta := i^\ast\beta\colon i^\ast Q\to i^\ast(G[n]) = (i^\ast G)[n]$.
    Recall that $i^\ast$ is exact and $i^\ast j_! = 0$.
    Applying $i^\ast$ to \eqref{eq:onecell_triangle} gives a distinguished triangle
    \begin{displaymath}
        0\to i^\ast Q\xrightarrow{\theta} i^\ast(G[n])\to 0[1],
    \end{displaymath}
    and so $\theta$ is an isomorphism.
    Moreover, applying $i^\ast$ to $\eta_Y\circ q = g\circ \beta$ and using the triangle identity $\overline{\varepsilon}_{i^\ast Y}\circ i^\ast(\eta_Y) = 1_{i^\ast Y}$, it follows that
    \begin{displaymath}
        i^\ast q
        = \overline{\varepsilon}_{i^\ast Y}\circ i^\ast(\eta_Y)\circ i^\ast q
        = \overline{\varepsilon}_{i^\ast Y}\circ i^\ast(g)\circ i^\ast\beta
        = f^\prime\circ \theta.
    \end{displaymath}

    Lastly, we check that $Q$ lies in $\mathcal{P}$.
    By \eqref{eq:onecell_triangle}, $Q\in j_! W\star G[n]$.
    Since $\mathcal{P}$ is closed under extensions, it suffices to show that $j_! W, G[n]\in \mathcal{P}$.
    The functor $j_!$ is exact, preserves small coproducts, and satisfies $j_! G_U\in \overline{\langle G\rangle}^{[0,0]}$ by \Cref{setup:standing}.
    By \Cref{lem:closures}\eqref{lem:closures3}, $j_! W\in \overline{\langle G\rangle}^{[1-3A,2A+1]}$.
    As $A\geq 1$, we have $2A+1\leq 3A = B$, so \Cref{lem:closures}\eqref{lem:closures2} gives $j_! W\in \overline{\langle G\rangle}^{[-B,B]}$.
    Moreover, by \Cref{lem:glued}, $W\in \mathcal{D}_U^{\leq K}$ implies $j_! W\in \mathcal{D}^{\leq K}$.
    Since $-A-1\leq n\leq A$, one has $-A\leq -n\leq A+1$, and so $G[n] = G[-(-n)]\in G[-A,A+1]$.
    Applying \Cref{lem:closures}\eqref{lem:closures2} with $A+1\leq 3A$, $G[n]\in \overline{\langle G\rangle}^{[-A,A+1]}\subseteq \overline{\langle G\rangle}^{[-B,B]}$.
    As $G\in \mathcal{D}^{\leq 0}$ by \Cref{setup:standing}, $G[n]\in \mathcal{D}^{\leq -n}\subseteq \mathcal{D}^{\leq A+1} = \mathcal{D}^{\leq K}$.
    Thus, $Q\in \mathcal{P}$, and so $(q,\theta)$ is a $\mathcal{P}$-lift of $f^\prime$.
\end{proof}

\begin{notation}
    Assume \Cref{setup:standing}. 
    Denote by $\mathcal{C}\subseteq \mathcal{D}_Z$ the full subcategory of objects $E^\prime$ such that every morphism $f^\prime \colon E^\prime \to i^\ast Y$ with $Y\in \mathcal{D}^{\leq K}$ admits a $\mathcal{P}$-lift.
\end{notation}

\begin{proposition}
    \label{prop:closure}
    Assume \Cref{setup:standing}. 
    The subcategory $\mathcal{C}$ is strictly full, closed under small coproducts and extensions, and contains $(i^\ast G)[n]$ for $-A-1\leq n\leq A$.
    Consequently,
    \begin{displaymath}
        \overline{\langle i^\ast G\rangle}^{[-A,A]}\subseteq \operatorname{Coprod}((i^\ast G)[-A,A+1])\subseteq \mathcal{C}.
    \end{displaymath}
\end{proposition}

\begin{proof}
    First we show that $\mathcal{C}$ is strictly full.
    Let $\sigma\colon E^\prime\to E^{\prime\prime}$ be an isomorphism with $E^{\prime\prime}\in \mathcal{C}$.
    Fix $f^\prime\colon E^\prime\to i^\ast Y$ with $Y\in \mathcal{D}^{\leq K}$.
    Since $E^{\prime\prime}\in \mathcal{C}$, $f^\prime\circ \sigma^{-1}$ has a $\mathcal{P}$-lift $(q,\theta)$.
    Then $\sigma^{-1}\circ \theta$ is an isomorphism, and $i^\ast q = f^\prime\circ (\sigma^{-1}\circ \theta)$.
    Thus, $(q,\sigma^{-1}\circ \theta)$ is a $\mathcal{P}$-lift of $f^\prime$. Hence, $E^\prime\in \mathcal{C}$.

    By \Cref{lem:onecell}, $(i^\ast G)[n] \in \mathcal{C}$  for $-A-1\leq n\leq A$.

    Next we show that $\mathcal{C}$ is closed under small coproducts.
    Choose a set $\Lambda$.
    Pick some $E^\prime_\lambda\in \mathcal{C}$ for each $\lambda\in \Lambda$.
    Put $E^\prime := \bigoplus_{\lambda\in \Lambda} E^\prime_\lambda$.
    Denote by $\iota_\lambda\colon E^\prime_\lambda\to E^\prime$ the inclusions.
    Fix $f^\prime\colon E^\prime\to i^\ast Y$ such that $Y\in \mathcal{D}^{\leq K}$.
    For each $\lambda \in \Lambda$, there exists a $\mathcal{P}$-lift $(q_\lambda\colon E_\lambda\to Y,\theta_\lambda)$ of $f^\prime\circ \iota_\lambda$.
    Define $E := \bigoplus_{\lambda\in \Lambda} E_\lambda$.
    Set $\kappa_\lambda\colon E_\lambda\to E$ to be the inclusions for each $\lambda \in \Lambda$.
    Since $\mathcal{P}$ is closed under small coproducts, $E\in \mathcal{P}$.
    Denote by $q\colon E\to Y$ the morphism with $q\circ \kappa_\lambda = q_\lambda$ for each $\lambda \in \Lambda$.
    As $i^\ast$ preserves small coproducts, $i^\ast E$ is the coproduct of the $i^\ast E_\lambda$ with inclusions $i^\ast(\kappa_\lambda)$ for $\lambda \in \Lambda$.
    Hence, there exists $\theta\colon i^\ast E\to E^\prime$ with $\theta\circ i^\ast(\kappa_\lambda) = \iota_\lambda\circ \theta_\lambda$ for each $\lambda \in \Lambda$.
    Note that $\theta$ is an isomorphism because it is a coproduct of the isomorphisms $\theta_\lambda$.
    For each $\lambda \in \Lambda$,
    \begin{displaymath}
        i^\ast q\circ i^\ast(\kappa_\lambda)
        = i^\ast(q_\lambda)
        = f^\prime\circ \iota_\lambda\circ \theta_\lambda
        = f^\prime\circ \theta\circ i^\ast(\kappa_\lambda).
    \end{displaymath}
    Hence, by the universal property of coproducts, $i^\ast q = f^\prime\circ \theta$.
    Thus, $(q,\theta)$ is a $\mathcal{P}$-lift of $f^\prime$.

    Now we prove that $\mathcal{C}$ is closed under extensions.
    Consider a distinguished triangle
    \begin{displaymath}
        E^\prime_1\xrightarrow{u^\prime} E^\prime\to E^\prime_2\to E^\prime_1[1]
    \end{displaymath}
    where $E^\prime_1,E^\prime_2\in \mathcal{C}$.
    Fix $f^\prime\colon E^\prime\to i^\ast Y$ with $Y\in \mathcal{D}^{\leq K}$.
    We will lift $f^\prime$ in two stages: first on $E^\prime_1$, and then on $E^\prime_2$, into the cone of the first lift.

    We start with $E^\prime_1$.
    As $E^\prime_1\in \mathcal{C}$, $f^\prime_1 := f^\prime\circ u^\prime$ has a $\mathcal{P}$-lift $(q_1\colon E_1\to Y,\theta_1)$.
    Complete to a distinguished triangle
    \begin{displaymath}
        E_1\xrightarrow{q_1} Y\xrightarrow{r} Y^\prime\xrightarrow{s} E_1[1].
    \end{displaymath}
    Since $E_1\in \mathcal{P} \subseteq \mathcal{D}^{\leq K}$, $E_1[1]\in \mathcal{D}^{\leq K-1}\subseteq \mathcal{D}^{\leq K}$.
    As aisles are closed under extensions, $Y^\prime\in Y\star E_1[1]\subseteq \mathcal{D}^{\leq K}$.
    %%NOTE: This is the key point: the cone $Y^\prime$ satisfies the same bound as $Y$.
    Applying $i^\ast$ to this distinguished triangle gives the top row of the diagram
    \begin{displaymath}
        \begin{tikzcd}[column sep=large]
            i^\ast E_1 \ar[r, "i^\ast q_1"] \ar[d, "\theta_1"'] & i^\ast Y \ar[r, "i^\ast r"] \ar[d, equal] & i^\ast Y^\prime \ar[r, "i^\ast s"] \ar[d, equal] & (i^\ast E_1)[1] \ar[d, "{\theta_1[1]}"] \\
            E^\prime_1 \ar[r, "f^\prime_1"] & i^\ast Y \ar[r, "i^\ast r"] & i^\ast Y^\prime \ar[r, "{\theta_1[1]\circ i^\ast s}"] & E^\prime_1[1].
        \end{tikzcd}
    \end{displaymath}
    Since $i^\ast q_1 = f^\prime_1\circ \theta_1$, the left square commutes.
    By construction, the other two squares commute.
    As all vertical maps are isomorphisms, the bottom row is isomorphic to a distinguished triangle, and hence is distinguished.

    Next we find the morphism that $f^\prime$ induces on $E^\prime_2$.
    Again, we complete to a distinguished triangle
    \begin{displaymath}
        E^\prime\xrightarrow{f^\prime} i^\ast Y\xrightarrow{p^\prime} R^\prime\to E^\prime[1].
    \end{displaymath}
    We apply \Cref{rmk:octahedral}\eqref{rmk:octahedraloc} to $E^\prime_1\xrightarrow{u^\prime} E^\prime\xrightarrow{f^\prime} i^\ast Y$.
    For this, we use the given distinguished triangle on $u^\prime$, the chosen distinguished triangle on $f^\prime$, and the bottom row of the diagram above as the distinguished triangle on $f^\prime\circ u^\prime = f^\prime_1$.
    This gives a distinguished triangle
    \begin{displaymath}
        E^\prime_2\xrightarrow{f^\prime_2} i^\ast Y^\prime\xrightarrow{h} R^\prime\to E^\prime_2[1]
    \end{displaymath}
    such that the following diagram commutes:
    \begin{displaymath}
        \begin{tikzcd}
            i^\ast Y \ar[r, "i^\ast r"] \ar[dr, "p^\prime"'] & i^\ast Y^\prime \ar[d, "h"] \\
            & R^\prime.
        \end{tikzcd}
    \end{displaymath}

    Now we lift $f^\prime_2$.
    Since $E^\prime_2\in \mathcal{C}$ and $Y^\prime\in \mathcal{D}^{\leq K}$, $f^\prime_2$ has a $\mathcal{P}$-lift $(q_2\colon E_2\to Y^\prime,\theta_2)$.
    Choose a distinguished triangle
    \begin{displaymath}
        E_2\xrightarrow{q_2} Y^\prime\xrightarrow{s_2} Y^{\prime\prime}\to E_2[1].
    \end{displaymath}

    Next we combine the two lifts.
    Complete to a distinguished triangle
    \begin{displaymath}
        E\xrightarrow{q} Y\xrightarrow{s_2\circ r} Y^{\prime\prime}\to E[1].
    \end{displaymath}
    The cocones of $r$, $s_2$, and $s_2\circ r$ are respectively $E_1$, $E_2$, and $E$.
    Hence, \Cref{rmk:octahedral}\eqref{rmk:octahedralof} gives a distinguished triangle $E_1\to E\to E_2\to E_1[1]$.
    Since $\mathcal{P}$ is closed under extensions, $E\in E_1\star E_2\subseteq \mathcal{P}$.

    It remains to compare $i^\ast q$ with $f^\prime$.
    We do this in two steps.
    First, applying $i^\ast$ to the distinguished triangle on $q_2$ gives the top row of the diagram
    \begin{displaymath}
        \begin{tikzcd}[column sep=large]
            i^\ast E_2 \ar[r, "i^\ast q_2"] \ar[d, "\theta_2"'] & i^\ast Y^\prime \ar[r, "i^\ast s_2"] \ar[d, equal] & i^\ast Y^{\prime\prime} \ar[r] \ar[d, dashed, "\chi"] & (i^\ast E_2)[1] \ar[d, "{\theta_2[1]}"] \\
            E^\prime_2 \ar[r, "f^\prime_2"] & i^\ast Y^\prime \ar[r, "h"] & R^\prime \ar[r] & E^\prime_2[1].
        \end{tikzcd}
    \end{displaymath}
    Note that the bottom row is the distinguished triangle on $f^\prime_2$ found above.
    As $i^\ast q_2 = f^\prime_2\circ \theta_2$, the left square commutes.
    By TR3, there exists a morphism $\chi\colon i^\ast Y^{\prime\prime}\to R^\prime$ making the diagram commute.
    Hence, $\chi\circ i^\ast s_2 = h$.
    Since $\theta_2$ and the identity are isomorphisms, so is  $\chi$.
    Combine with $h\circ i^\ast r = p^\prime$ to obtain
    \begin{displaymath}
        \chi\circ i^\ast(s_2\circ r) = h\circ i^\ast r = p^\prime.
    \end{displaymath}
    Applying $i^\ast$ to the distinguished triangle on $q$ gives 
    \begin{displaymath}
        \begin{tikzcd}[column sep=large]
            i^\ast E \ar[r, "i^\ast q"] \ar[d, dashed, "\theta"'] & i^\ast Y \ar[r, "i^\ast(s_2\circ r)"] \ar[d, equal] & i^\ast Y^{\prime\prime} \ar[r] \ar[d, "\chi"] & (i^\ast E)[1] \ar[d, dashed, "{\theta[1]}"] \\
            E^\prime \ar[r, "f^\prime"] & i^\ast Y \ar[r, "p^\prime"] & R^\prime \ar[r] & E^\prime[1].
        \end{tikzcd}
    \end{displaymath}
    The bottom row is the distinguished triangle on $f^\prime$, whereas the middle square commutes by the diagram above.
    By TR3, applied after rotating both rows, there exists a morphism $\theta\colon i^\ast E\to E^\prime$ making the diagram commute.
    Hence, $f^\prime\circ \theta = i^\ast q$.
    As before, $\theta$ is an isomorphism, because $1$ and $\chi$ are.
    Since $E\in \mathcal{P}$, $(q,\theta)$ is a $\mathcal{P}$-lift of $f^\prime$.
    Hence, $E^\prime\in \mathcal{C}$.

    Finally, we prove the two inclusions.
    The objects of $(i^\ast G)[-A,A+1]$ are $(i^\ast G)[-k]$ with $-A\leq k\leq A+1$.
    For $n = -k$, these are the objects $(i^\ast G)[n]$ with $-A-1\leq n\leq A$.
    Hence, by \Cref{lem:onecell}, these lie in $\mathcal{C}$.
    Moreover, by definition, $\operatorname{Coprod}((i^\ast G)[-A,A+1])$ is the smallest strictly full subcategory containing these objects that is closed under small coproducts and extensions.
    Since $\mathcal{C}$ has all these properties, $\operatorname{Coprod}((i^\ast G)[-A,A+1])\subseteq \mathcal{C}$.
    The first inclusion is \Cref{lem:closures}\eqref{lem:closures5}.
\end{proof}

\begin{remark}
    \Cref{prop:closure} replaces the induction on the number of cones in \cite[Lemmas 5.4 and 5.5]{Burke/Neeman/Pauwels:2023}, which requires approximability.
\end{remark}

%%%%%%%%%%%%%%%%%%%%%%%%%%%%%%%%%%%%%%%%%%%
\section{Proofs of results}
\label{sec:proof_of_results}
%%%%%%%%%%%%%%%%%%%%%%%%%%%%%%%%%%%%%%%%%%%

\begin{proof}
    [Proof of \Cref{thm:glueing_weak_approx}]
    First, let $\mathcal{D}$ be weakly approximable.
    Denote by $\mathcal{D}^{\leq 0}$ the aisle compactly generated by $G^\prime$.
    Note that $G^\prime\in \mathcal{D}^{\leq 0}$.
    By \cite[Lemma 3.7]{Burke/Neeman/Pauwels:2023}, there exists an integer $A^\prime>0$ such that $(G^\prime,\mathcal{D}^{\leq 0},A^\prime)$ satisfies \Cref{def:weak_approx}\eqref{def:weak_approx1}.
    It follows that $\operatorname{Hom}(G^\prime,\mathcal{D}^{\leq -A^\prime})=0$.
    For $n\geq A^\prime$, $G^\prime[n]\in \mathcal{D}^{\leq -n}\subseteq \mathcal{D}^{\leq -A^\prime}$.
    Hence, for all $n\geq A^\prime$, $\operatorname{Hom}(G^\prime,G^\prime[n])=0$.
    This argument applies to every compact generator of $\mathcal{D}$.

    Conversely, let $\operatorname{Hom}(G^\prime,G^\prime[n])=0$ for $n\gg 0$.
    Hence, we are in the setting of \Cref{setup:standing}, and we use its notation. 
    By \Cref{lem:normalize}\eqref{lem:normalize2}, $G$ is a compact generator of $\mathcal{D}$.
    We show that $(G,\mathcal{D}^{\leq 0},3A)$ is weak approximation data for $\mathcal{D}$.

    We start with \Cref{def:weak_approx}\eqref{def:weak_approx1}.
    Since $G\in \mathcal{D}^{\leq 0}$, $G\in \mathcal{D}^{\leq 3A}$.
    As $\mathcal{D}^{\leq -3A}\subseteq \mathcal{D}^{\leq -A}$, \Cref{setup:standing}\eqref{setup:standinga3} implies $\operatorname{Hom}(G,\mathcal{D}^{\leq -3A})=0$.

    Now we verify \Cref{def:weak_approx}\eqref{def:weak_approx2}.
    Fix $X\in \mathcal{D}^{\leq 0}$.
    We will approximate $X$ in two stages: first its image $i^\ast X$ in $\mathcal{D}_Z$, and then what remains in $\mathcal{D}_U$.

    We start on $\mathcal{D}_Z$. 
    By \Cref{lem:glued}, $i^\ast X\in \mathcal{D}_Z^{\leq 0}$.
    By \Cref{setup:standing}\eqref{setup:standinga1}, there exists a distinguished triangle
    \begin{displaymath}
        E^\prime\xrightarrow{f^\prime} i^\ast X\to D^\prime\to E^\prime[1]
    \end{displaymath}
    where $E^\prime\in \overline{\langle i^\ast G\rangle}^{[-A,A]}$ and $D^\prime\in \mathcal{D}_Z^{\leq -1}$.
    By \Cref{prop:closure}, $E^\prime\in \mathcal{C}$.
    Since $X\in \mathcal{D}^{\leq 0}\subseteq \mathcal{D}^{\leq K}$, $f^\prime$ has a $\mathcal{P}$-lift $(q_1\colon E_1\to X,\theta)$.
    Complete to a distinguished triangle
    \begin{displaymath}
        E_1\xrightarrow{q_1} X\xrightarrow{r} N\to E_1[1].
    \end{displaymath}
    Since $E_1\in \mathcal{P}\subseteq \mathcal{D}^{\leq K}$, $E_1[1]\in \mathcal{D}^{\leq K-1}\subseteq \mathcal{D}^{\leq K}$.
    As aisles are closed under extensions, $N\in X\star E_1[1]\subseteq \mathcal{D}^{\leq K}$.
    Applying $i^\ast$ to the distinguished triangle on $q_1$ gives the top row of the diagram
    \begin{displaymath}
        \begin{tikzcd}[column sep=large]
            i^\ast E_1 \ar[r, "i^\ast q_1"] \ar[d, "\theta"'] & i^\ast X \ar[r, "i^\ast r"] \ar[d, equal] & i^\ast N \ar[r] \ar[d, dashed, "\mu"] & (i^\ast E_1)[1] \ar[d, "{\theta[1]}"] \\
            E^\prime \ar[r, "f^\prime"] & i^\ast X \ar[r] & D^\prime \ar[r] & E^\prime[1],
        \end{tikzcd}
    \end{displaymath}
    whose bottom row is the distinguished triangle on $f^\prime$.
    The left square commutes because $i^\ast q_1 = f^\prime\circ \theta$.
    By TR3, there exists $\mu\colon i^\ast N\to D^\prime$ making the diagram commute.
    Since $\theta$ and the identity are isomorphisms, so is $\mu$.
    By definition, aisles are strictly full, and so $i^\ast N\in \mathcal{D}_Z^{\leq -1}$.

    Next we treat the case on $\mathcal{D}_U$.
    By \Cref{lem:glued}, $j^\ast N\in \mathcal{D}_U^{\leq K}$, and so $(j^\ast N)[K]\in \mathcal{D}_U^{\leq 0}$.
    Apply \Cref{lem:iterate} to the weak approximation data $(G_U,\mathcal{D}_U^{\leq 0},A)$ of \Cref{setup:standing}\eqref{setup:standinga2} and the object $(j^\ast N)[K]$, with $K+1$ in place of $m$.
    This gives a distinguished triangle
    \begin{displaymath}
        V^\prime\to (j^\ast N)[K]\to R^\prime\to V^\prime[1]
    \end{displaymath}
    where $V^\prime\in \overline{\langle G_U\rangle}^{[-K-A,A]}$ and $R^\prime\in \mathcal{D}_U^{\leq -K-1}$.
    Define $V := V^\prime[-K]$ and $R := R^\prime[-K]$.
    Shifting by $[-K]$ yields a distinguished triangle
    \begin{displaymath}
        V\xrightarrow{c} j^\ast N\to R\to V[1].
    \end{displaymath}
    By \Cref{lem:closures}\eqref{lem:closures2},
    \begin{displaymath}
        V\in \overline{\langle G_U\rangle}^{[-A,A+K]} = \overline{\langle G_U\rangle}^{[-A,2A+1]}.
    \end{displaymath}
    Moreover, $R\in \mathcal{D}_U^{\leq -K-1}[-K] = \mathcal{D}_U^{\leq -1}$.

    Consider the composite
    \begin{displaymath}
        \begin{tikzcd}
            j_! V \ar[r, "j_!(c)"] \ar[rr, bend left=30, "\psi"] & j_! j^\ast N \ar[r, "\epsilon_N"] & N.
        \end{tikzcd}
    \end{displaymath}
    Complete to a distinguished triangle
    \begin{displaymath}
        j_! V\xrightarrow{\psi} N \xrightarrow{\rho} D^{\prime\prime}\to j_! V[1].
    \end{displaymath}
    We claim that $D^{\prime\prime}\in \mathcal{D}^{\leq -1}$.
    Since $j_!$ is exact, applying it to the distinguished triangle on $c$ gives a distinguished triangle
    \begin{displaymath}
        j_! V\xrightarrow{j_!(c)} j_! j^\ast N\to j_! R\to j_! V[1].
    \end{displaymath}
    Hence, $\operatorname{cone}(j_!(c)) \cong j_! R$.
    By \eqref{eq:gluing_triangles}, $\operatorname{cone}(\epsilon_N)\cong i_\ast i^\ast N$.
    Applying \Cref{rmk:octahedral}\eqref{rmk:octahedraloc} to the composite above gives a distinguished triangle
    \begin{displaymath}
        j_! R\to D^{\prime\prime}\to i_\ast i^\ast N\to j_! R[1].
    \end{displaymath}
    By \Cref{lem:glued}, $R\in \mathcal{D}_U^{\leq -1}$ implies $j_! R\in \mathcal{D}^{\leq -1}$.
    Similarly, $i^\ast N\in \mathcal{D}_Z^{\leq -1}$ implies $i_\ast i^\ast N \in \mathcal{D}^{\leq -1}$.
    As aisles are closed under extensions, $D^{\prime\prime}\in j_! R\star i_\ast i^\ast N \subseteq \mathcal{D}^{\leq -1}$.

    We now assemble the two stages.
    Consider the composite
    \begin{displaymath}
        \begin{tikzcd}
            X \ar[r, "r"] \ar[rr, bend left=30, "\rho\circ r"] & N \ar[r, "\rho"] & D^{\prime\prime}.
        \end{tikzcd}
    \end{displaymath}
    Complete to a distinguished triangle
    \begin{displaymath}
        E\to X\xrightarrow{\rho\circ r} D^{\prime\prime}\to E[1].
    \end{displaymath}
    The cocones of $r$, $\rho$, and $\rho\circ r$ are respectively $E_1$, $j_! V$, and $E$.
    Hence, \Cref{rmk:octahedral}\eqref{rmk:octahedralof} gives a distinguished triangle $E_1\to E\to j_! V\to E_1[1]$.
    Note that $E\in E_1\star j_! V$.

    Lastly we bound the range of $E$.
    Since $E_1\in \mathcal{P}$ and $B = 3A$, $E_1\in \overline{\langle G\rangle}^{[-3A,3A]}$.
    Note that $j_!$ is exact and preserves small coproducts.
    By \Cref{setup:standing}, $j_! G_U\in \overline{\langle G\rangle}^{[0,0]}$.
    By \Cref{lem:closures}\eqref{lem:closures3}, $j_! V\in \overline{\langle G\rangle}^{[-A,2A+1]}$.
    As $2A+1\leq 3A$, \Cref{lem:closures}\eqref{lem:closures2} gives $j_! V\in \overline{\langle G\rangle}^{[-3A,3A]}$.
    By \Cref{lem:closures}\eqref{lem:closures1}, $\overline{\langle G\rangle}^{[-3A,3A]}$ is closed under extensions, and so $E\in \overline{\langle G\rangle}^{[-3A,3A]}$.

    Thus, $E\to X\to D^{\prime\prime}\to E[1]$ is a distinguished triangle with $E\in \overline{\langle G\rangle}^{[-3A,3A]}$ and $D^{\prime\prime}\in \mathcal{D}^{\leq -1}$.
    As $X\in \mathcal{D}^{\leq 0}$ was arbitrary, $(G,\mathcal{D}^{\leq 0},3A)$ is weak approximation data for $\mathcal{D}$.
    Therefore, $\mathcal{D}$ is weakly approximable.
\end{proof}

\begin{remark}
    The proof of \Cref{thm:glueing_weak_approx} bootstraps that of \cite[Theorem 5.1]{Burke/Neeman/Pauwels:2023}.
    The `only if' direction is \cite[Remark 3]{Burke/Neeman/Pauwels:2023}.
    For the `if' direction, the two stages compare as follows.
    \begin{itemize}
        \item (On $\mathcal{D}_Z$) In \cite{Burke/Neeman/Pauwels:2023}, the approximation $E^\prime \to i^\ast X$ has $E^\prime \in \overline{\langle i^\ast G\rangle}^{[-A,A]}_A$ (see loc.\ cit.\ for notation). 
        By \cite[Lemma 5.5]{Burke/Neeman/Pauwels:2023}, it can be lifted.
        In fact, loc.\ cit.\ proceeds by induction on the number of cones, which requires approximability.
        Here $E^\prime$ is only known to lie in $\overline{\langle i^\ast G\rangle}^{[-A,A]}$.
        Instead, it is lifted by \Cref{prop:closure}.
        \item (On $\mathcal{D}_U$) The octahedra are those in the last paragraph of the proof of \cite[Theorem 5.1]{Burke/Neeman/Pauwels:2023}.
        Here \Cref{lem:iterate} is in place of the estimate for approximable categories used in loc.\ cit.\
    \end{itemize}
\end{remark}

\begin{remark}
    Let $X$ be a quasi-compact quasi-separated algebraic space. 
    Pick a closed subset $S\subseteq |X|$ with quasi-compact complement. 
    Consider the following subcategories
    \begin{itemize}
        \item $D^{\leq 0}(X)$ consists of objects $E\in D(X)$ such that $\mathcal{H}^j (E)=0$ for $j >0$
        \item $D^{\geq 0}(X)$ consists of objects $E\in D(X)$ such that $\mathcal{H}^j (E)=0$ for $j <0$. 
    \end{itemize}
    We call $(D^{\leq 0}(X), D^{\geq 0}(X))$ the \textbf{standard $t$-structure} on $D(X)$. 
    The standard truncations of an object $E\in D_{\operatorname{qc}}(X)$ have quasi-coherent cohomology.
    Moreover, the support is defined via cohomology sheaves, so the truncations of an object supported on $S$ are supported on $S$.
    See `canonical truncations' in \cite[\href{https://stacks.math.columbia.edu/tag/0118}{Tag 0118}]{StacksProject}.
    Hence, the standard $t$-structure induces a $t$-structure on $D_{\operatorname{qc},S}(X)$, which is denoted by $(D^{\leq 0}_{\operatorname{qc},S}(X), D^{\geq 0}_{\operatorname{qc},S}(X))$.
    Write $D^?_{\operatorname{qc},S}(X)$ for $(D_{\operatorname{qc},S}(X))^?$ in \Cref{sec:prelim_t-structures_t_structures} with $?\in \{-,+,b\}$. 
\end{remark}

\begin{definition}
    \label{def:wa_spaces}
    Let $X$ be a quasi-compact quasi-separated algebraic space. 
    We say that $X$ satisfies \textbf{W.A.}\ if $D_{\operatorname{qc},S}(X)$ is weakly approximable for every closed subset $S\subseteq |X|$ with quasi-compact complement.
    We say $X$ satisfies \textbf{standard W.A.} if for every closed subset $S\subseteq |X|$ with quasi-compact complement, there exists weak approximation data $(G,D^{\leq 0}_{\operatorname{qc},S}(X),A)$.
\end{definition}

\begin{remark}
    It is easy to see that standard W.A.\ implies W.A.\ in the sense of \Cref{def:wa_spaces}.
\end{remark}

\begin{lemma}
\label{lem:compact_ext_vanishing}
	Let $X$ be a quasi-compact and quasi-separated algebraic space.
	If $G$ is a compact object of $D_{\operatorname{qc}}(X)$, then
	there exists an integer $N$ such that $\operatorname{Hom}(G,G[n])=0$ for every $n>N$.
\end{lemma}

\begin{proof}
	By \cite[\href{https://stacks.math.columbia.edu/tag/09M8}{Tag 09M8}]{StacksProject}, $G$ is perfect. 
    Moreover, \cite[\href{https://stacks.math.columbia.edu/tag/08J7}{Tag 08J7}]{StacksProject} yields 
    \begin{displaymath}
        \operatorname{Hom}(G,G[n])
        \cong \operatorname{Hom}(\mathcal{O}_X, \operatorname{\mathbb{R}\mathcal{H}\! \mathit{om}}(G,G[n])).
    \end{displaymath}
    Note that $\operatorname{\mathbb{R}\mathcal{H}\! \mathit{om}}(G,G[n])$ is perfect \cite[\href{https://stacks.math.columbia.edu/tag/08JJ}{Tag 08JJ}]{StacksProject}.
    Applying \cite[\href{https://stacks.math.columbia.edu/tag/0GFH}{Tags 0GFH} and \href{https://stacks.math.columbia.edu/tag/0GFI}{0GFI}]{StacksProject}, it follows that 
    \begin{displaymath}
        \operatorname{Hom}(\mathcal{O}_X, \operatorname{\mathbb{R}\mathcal{H}\! \mathit{om}}(G,G[n])) \cong 0.
    \end{displaymath}
    for $|n| \gg 0$.
    This completes the proof.
\end{proof}

\begin{lemma}
    \label{lem:perf_vanishing}
    Let $X$ be a quasi-compact quasi-separated algebraic space. 
    Choose $S\subseteq |X|$ closed with quasi-compact complement.
    If $G^\prime$ is a compact generator of $D_{\operatorname{qc},S}(X)$, then $\operatorname{Hom}(G^\prime,G^\prime[n])=0$ for $n\gg 0$.
\end{lemma}

\begin{proof}
    This follows from \cite[Remark A.5]{GuisadoVillaalgordo/Lank:2026}, \cite[\href{https://stacks.math.columbia.edu/tag/09M8}{Tag 09M8}]{StacksProject}, and \Cref{lem:compact_ext_vanishing}.
\end{proof}

\begin{lemma}
	\label{lem:qcqs_open_pushforward}
	Let $j\colon U\to X$ be a quasi-compact open immersion between quasi-compact quasi-separated algebraic spaces. Then
	$\mathbf{R}j_\ast\colon D_{\operatorname{qc}}(U)\to
	D_{\operatorname{qc}}(X)$ preserves small coproducts.
	There is an integer $r\geq 0$ such that
	\begin{displaymath}
		\mathbf{R}j_\ast D^{\leq 0}_{\operatorname{qc}}(U)
		\subseteq D^{\leq r}_{\operatorname{qc}}(X).
	\end{displaymath}
	For any quasi-compact open subset $V\subseteq |U|$, the same
	statements hold for the restricted functor
	$\mathbf{R}j_\ast\colon D_{\operatorname{qc},|U|\setminus V}(U)\to
	D_{\operatorname{qc},|X|\setminus V}(X)$.
\end{lemma}

\begin{proof}
    By \cite[\href{https://stacks.math.columbia.edu/tag/08FB}{Tag 08FB}]{StacksProject}, $\mathbf{R}j_\ast$ preserves small coproducts.
    Let $s\colon V \to X$ be an \'{e}tale presentation by an affine scheme.
    Consider the fibered square 
    \begin{displaymath}
        % https://q.uiver.app/#q=WzAsNCxbMSwwLCJWIl0sWzEsMSwiXFxtYXRoY2Fse1h9LiJdLFswLDEsIlxcbWF0aGNhbHtVfSJdLFswLDAsIlxcbWF0aGNhbHtVfVxcdGltZXNfe1xcbWF0aGNhbHtYfX0gViJdLFswLDEsInMiXSxbMiwxLCJqIiwyXSxbMywwLCJqXlxccHJpbWUiXSxbMywyLCJzXlxccHJpbWUiLDJdXQ==
        \begin{tikzcd}
            {U\times_{X} V} & V \\
            {U} & {X.}
            \arrow["{j^\prime}", from=1-1, to=1-2]
            \arrow["{s^\prime}"', from=1-1, to=2-1]
            \arrow["s", from=1-2, to=2-2]
            \arrow["j"', from=2-1, to=2-2]
        \end{tikzcd}
    \end{displaymath}
    As $j$ is representable by schemes, $U\times_{X} V$ is a scheme.
    Base change implies $j^\prime$ is an open immersion, and that $s^\prime$ is a smooth surjection.
    By \cite[Proposition 3.9.2]{Lipman:2009}, there is an integer $r\geq 0$ such that
	\begin{displaymath}
		\mathbf{R}j_\ast D^{\leq 0}_{\operatorname{qc}}(U\times_{X} V)
		\subseteq D^{\leq r}_{\operatorname{qc}}(V).
	\end{displaymath}
    Pick $E\in D^{\leq 0}_{\operatorname{qc}}(U)$.
    By \cite[\href{https://stacks.math.columbia.edu/tag/08IR}{Tag 08IR}]{StacksProject}, $\mathbf{L}s^\ast \mathbf{R}j_\ast E \cong \mathbf{R}j^\prime_\ast \mathbf{L}(s^\prime)^\ast E$.
    It follows that $\mathbf{R}j^\prime_\ast \mathbf{L}(s^\prime)^\ast E \in D^{\leq r}_{\operatorname{qc}}(V)$ implies $\mathbf{R}j_\ast E \in D^{\leq r}_{\operatorname{qc}}(X)$ \cite[Lemma 2.15]{GuisadoVillaalgordo/Lank:2026}.
    This proves the first claim, whereas the last follows from supports.
\end{proof}

\begin{proposition}
    \label{prop:A8}
    Consider an elementary distinguished square $j\colon U\to X$ and $f\colon Y\to X$ of quasi-compact quasi-separated algebraic spaces (here $j$ is an open immersion, see \cite[\href{https://stacks.math.columbia.edu/tag/08GM}{Tag 08GM}]{StacksProject}) 
    If $U$ and $Y$ satisfy (resp.\ standard) W.A., then so does $X$.
\end{proposition}

\begin{proof}
    Fix a closed subset $S\subseteq |X|$ with quasi-compact complement.
    Set $T := |X|\setminus |U|$.
    By \cite[Lemma A.4]{GuisadoVillaalgordo/Lank:2026}, there exists a recollement
    \begin{displaymath}
        \begin{tikzcd}[column sep=5.5em]
            {D_{\operatorname{qc},|U|\cap S}(U)} & {D_{\operatorname{qc},S}(X)} & {D_{\operatorname{qc},S\cap T}(X)}
            \arrow["{\mathbf{R}j_\ast}" description, from=1-1, to=1-2]
            \arrow["{\mathbf{L}j^\ast}"', bend right=35, from=1-2, to=1-1]
            \arrow["{j^!_S}", bend left=35, from=1-2, to=1-1]
            \arrow["{i^!}" description, from=1-2, to=1-3]
            \arrow["{i_\ast}"', bend right=35, from=1-3, to=1-2]
            \arrow["{i^\ast}", bend left=35, from=1-3, to=1-2]
        \end{tikzcd}
    \end{displaymath}
    where $i_\ast$ is the natural inclusion.
    We apply \Cref{thm:glueing_weak_approx} to this recollement.
    In the notation of \Cref{setup:recollement}, the roles are as follows:
    \begin{itemize}
        \item $\mathcal{D}_Z = D_{\operatorname{qc},|U|\cap S}(U)$, where $\mathbf{L}j^\ast$, $\mathbf{R}j_\ast$, $j^!_S$ (of the diagram above) are respectively identified with $i^\ast$, $i_\ast$, $i^!$
        \item $\mathcal{D} = D_{\operatorname{qc},S}(X)$
        \item $\mathcal{D}_U = D_{\operatorname{qc},S\cap T}(X)$, where $i_\ast$, $i^!$, $i^\ast$ are respectively identified with $j_!$, $j^\ast$, $j_\ast$ (in \Cref{setup:recollement}).
    \end{itemize}

    First, $|U|\setminus (|U|\cap S) = |U|\cap (|X|\setminus S)$.
    Since $|U|\cap (|X|\setminus S)$ is the intersection of two quasi-compact opens of the quasi-separated topological space $|X|$, it is quasi-compact.
    As $U$ satisfies W.A., $\mathcal{D}_Z$ is weakly approximable.

    Next, apply \cite[Lemma A.6]{GuisadoVillaalgordo/Lank:2026} to $T\cap S$.
    This gives a triangulated equivalence
    \begin{displaymath}
        \mathbf{L}f^\ast\colon D_{\operatorname{qc},T\cap S}(X)\to D_{\operatorname{qc},f^{-1}(T\cap S)}(Y).
    \end{displaymath}
    Note that $|Y|\setminus f^{-1}(T\cap S) = f^{-1}(|U|\cup (|X|\setminus S))$.
    Since $Y$ is quasi-compact and $X$ is quasi-separated, $f$ is quasi-compact \cite[\href{https://stacks.math.columbia.edu/tag/073B}{Tag 073B}]{StacksProject}.
    As $|U|\cup (|X|\setminus S)$ is a quasi-compact open, $f^{-1}(|U|\cup (|X|\setminus S))$ is quasi-compact.
    Moreover, as $Y$ satisfies W.A., $D_{\operatorname{qc},f^{-1}(T\cap S)}(Y)$ is weakly approximable.
    By \cite[Remark 3.4(i)]{Canonaco/Neeman/Stellari:2024}, weak approximability is invariant under triangulated equivalence.
    Thus, $\mathcal{D}_U$ is weakly approximable.

    Observe that $|X|\setminus (S\cap T) = (|X|\setminus S)\cup |U|$, which is quasi-compact.
    Moreover, by \cite[\href{https://stacks.math.columbia.edu/tag/0AEC}{Tag 0AEC}]{StacksProject}, the three triangulated categories in the recollement are singly compactly generated.
    If $G^\prime$ is a compact generator of $D_{\operatorname{qc},S}(X)$, then \Cref{lem:perf_vanishing} says that $\operatorname{Hom}(G^\prime,G^\prime[n])=0$ for $n\gg 0$.
    Consequently, \Cref{thm:glueing_weak_approx} asserts that $D_{\operatorname{qc},S}(X)$ is weakly approximable.

    Suppose that $U$ and $Y$ satisfy standard W.A.
    By \cite[Facts 1.26(i)]{Neeman:2026}, the standard $t$-structures on the two outside terms lie in their respective preferred equivalence classes. For the term supported on $T\cap S$, use the $t$-exact equivalence induced by $\mathbf{L}f^\ast$.
    Denote by $\mathcal{D}^{\leq 0}_{\mathrm{gl}}$ the aisle obtained by gluing these two standard $t$-structures.

    In the proof of \Cref{thm:glueing_weak_approx}, the $t$-structure
    used on $D_{\operatorname{qc},S}(X)$ is glued from $t$-structures
    in the preferred equivalence classes of the two outside terms.
    See \Cref{setup:standing}. 
    Replacing either outside $t$-structure by an equivalent one changes the glued $t$-structure only up to equivalence. 
    Indeed, this follows from the description of its aisle in \Cref{lem:glued}, with the maximum of the two bounds. 
    %%NOTE: Shift bounds
    Hence, $\mathcal{D}^{\leq 0}_{\mathrm{gl}}$ is in the preferred equivalence class of $D_{\operatorname{qc},S}(X)$.

    We compare this glued $t$-structure with the standard one.
    By \Cref{lem:qcqs_open_pushforward}, there exists $r\geq 0$ such that $\mathbf{R}j_\ast D^{\leq 0}_{\operatorname{qc},|U|\cap S}(U) \subseteq D^{\leq r}_{\operatorname{qc},S}(X)$.
    If $E\in\mathcal{D}^{\leq 0}_{\mathrm{gl}}$, the recollement triangle
    \begin{displaymath}
        i_\ast i^! E\to E\to
        \mathbf{R}j_\ast\mathbf{L}j^\ast E\to (i_\ast i^! E)[1]
    \end{displaymath}
    implies $E\in D^{\leq r}_{\operatorname{qc},S}(X)$.
    Conversely, if $E\in D^{\leq 0}_{\operatorname{qc},S}(X)$, the same distinguished triangle implies $i^! E\in D^{\leq r+1}_{\operatorname{qc},S\cap T}(X)$.
    As $\mathbf{L}j^\ast E\in
    D^{\leq 0}_{\operatorname{qc},|U|\cap S}(U)$, it follows that
    \begin{displaymath}
        D^{\leq 0}_{\operatorname{qc},S}(X)
        \subseteq \mathcal{D}^{\leq r+1}_{\mathrm{gl}}
        \subseteq D^{\leq 2r+1}_{\operatorname{qc},S}(X).
    \end{displaymath}
    Hence, the glued and standard $t$-structures are equivalent.
    Since the standard $t$-structure belongs to the preferred class,
    \cite[Lemma 3.7]{Burke/Neeman/Pauwels:2023} gives the desired weak approximation data. 
    Therefore, $X$ satisfies standard W.A.\
    Since $S$ was arbitrary, the first assertion gives W.A.\ and the other argument gives standard W.A.\
\end{proof}

\begin{lemma}
    \label{lem:characterize_T-_c}
    Let $X$ be a quasi-compact quasi-separated algebraic space that satisfies standard W.A.
    Fix a closed subset $S\subseteq|X|$ with quasi-compact complement.
    Set $\mathcal{T}:=D_{\operatorname{qc},S}(X)$.
    Then $\mathcal{T}^-_c$ consists of the objects of $\mathcal{T}$ that are pseudocoherent in $D(X)$ and supported on $S$.
    If $X$ is Noetherian, $\mathcal{T}^?_c = D^?_{\operatorname{coh},S}(X)$ for $?\in \{-,b\}$.
\end{lemma}

\begin{proof}
    Since the standard $t$-structure appears in weak approximation data, \cite[Facts 1.26(i)]{Neeman:2026} implies it lies in the preferred equivalence class.
    Hence, we compute $\mathcal{T}^-_c$ and $\mathcal{T}^b_c$ with respect to it.
 
    Pick $E\in\mathcal{T}^-_c$ and $n>0$. 
    Choose a distinguished triangle
    \begin{displaymath}
        D\to E\to F\to D[1]
    \end{displaymath}
    with $D\in\mathcal{T}^c$ and $F\in\mathcal{T}^{\leq -n}$.
    By By \cite[\href{https://stacks.math.columbia.edu/tag/09M8}{Tag 09M8}]{StacksProject}, $G$ is perfect. , $D$ is perfect, and hence, pseudocoherent \cite[\href{https://stacks.math.columbia.edu/tag/08G8}{Tag 08G8}]{StacksProject}.
    Since $\mathcal{H}^j(F)=0$ for $j\geq 1-n$, $F$ is $(1-n)$-pseudocoherent.
    Hence, $E$ is $(1-n)$-pseudocoherent \cite[\href{https://stacks.math.columbia.edu/tag/08FV}{Tag 08FV}]{StacksProject}.
    As $n$ was arbitrary, $E$ is pseudocoherent.

    Conversely, let $E\in\mathcal{T}$ be pseudocoherent.
    Fix $n>0$.
    The open immersion with image $|X|\setminus S$ is quasi-compact \cite[\href{https://stacks.math.columbia.edu/tag/073B}{Tag 073B}]{StacksProject}.
    By \cite[\href{https://stacks.math.columbia.edu/tag/08HP}{Tag 08HP}]{StacksProject}, there exist a compact object $P$ supported on $S$ and $P\to E$ inducing isomorphisms on $\mathcal{H}^i$ for $i>1-n$ and a surjection on $\mathcal{H}^{1-n}$.
    Note that the cone of $P\to E$ lies in $\mathcal{T}^{\leq -n}$.
    Clearly, $P\in\mathcal{T}^c$.
    Therefore, $E\in\mathcal{T}^-_c$, which gives the first claim.

    We show the second claim. 
    By \cite[\href{https://stacks.math.columbia.edu/tag/08IK}{Tag 08IK}]{StacksProject}, the pseudocoherent objects of $D_{\operatorname{qc}}(X)$ are exactly those of $D^-_{\operatorname{coh}}(X)$.
    Indeed, this can be checked by pulling back on a smooth presentation.
    Hence, $\mathcal{T}^-_c = D^-_{\operatorname{coh}}(X) \cap \mathcal{T} = D^-_{\operatorname{coh},S}(X)$, and $\mathcal{T}^b_c = \mathcal{T}^b \cap\mathcal{T}^-_c =D^b_{\operatorname{coh},S}(X)$. 
    This completes the proof.
\end{proof}

\begin{proof}
    [Proof of \Cref{cor:weak_approx_spaces}]
    We establish standard W.A., which is stronger than the first claim.
    Fix $X$ as in the statement.
    Let $P$ be the property that an algebraic space satisfies standard W.A.
    We apply the induction principle \cite[\href{https://stacks.math.columbia.edu/tag/09IT}{Tag 09IT}]{StacksProject}.

    Set $W$ to be the schematic locus of $X$.
    By \cite[\href{https://stacks.math.columbia.edu/tag/06NH}{Tag 06N}]{StacksProject}, $W$ is a dense open subspace of $X$ which is a scheme.
    By \cite[Theorem 3.2(iii,iv)]{Neeman:2024}, $W$ satisfies standard W.A.

    Consider an elementary distinguished square $(W_1 \subseteq W_2,f\colon V \to W_2)$ such that each $W_i$ is a quasi-compact open subspace of $X$, $W\subseteq W_1$, $V$ is affine, and $P$ holds for $W_1$. 
    By \cite[Theorem 3.2(iii,iv)]{Neeman:2024}, $V$ satisfies standard W.A.
    Thus, \Cref{prop:A8} implies $W_2$ satisfies standard W.A.
    Consequently, $X$ satisfies standard W.A.
    Hence, the first claim follows, and the case $S=|X|$ gives the second.

    We prove the last claim.
    Here $X$ is Noetherian. 
    Fix a closed subset $S\subseteq|X|$ with quasi-compact complement.
    Set $\mathcal{T}:=D_{\operatorname{qc},S}(X)$.
    By \Cref{lem:characterize_T-_c}, $\mathcal{T}^-_c = D^-_{\operatorname{coh},S}(X)$ and $\mathcal{T}^b_c = D^b_{\operatorname{coh},S}(X)$. 
    The standard $t$-structure restricts to $D^-_{\operatorname{coh},S}(X)$. 
    Thus, we finish by \cite[Example 10.1.2]{Canonaco/Haesemeyer/Neeman/Stellari:2024} and the first claim.
    Indeed, by loc.\ cit., a weakly approximable category is coherent if some $t$-structure in the preferred class restricts to $\mathcal{T}^-_c$.
\end{proof}

\bibliographystyle{alpha}
\bibliography{mainbib}

\end{document}